%%%%%%%%%%%%%%%%%%%%%%%%%%%%%%%%%%%%%%%%%%%%%%%%%%%%%%%%%%%%
% 13.1 MATHEMATICAL MODELING
% The Composition of Advanced Mathematics
%%%%%%%%%%%%%%%%%%%%%%%%%%%%%%%%%%%%%%%%%%%%%%%%%%%%%%%%%%%%

\section{Mathematical Modeling}
\label{sec:mathematical-modeling}

\prerequisite{
Functions, algebra, calculus, differential equations, probability,
statistics, optimization, numerical methods, dimensional analysis,
and familiarity with mathematical notation and scientific reasoning.
}

\learningobjectives{
By the end of this section, the reader should be able to:
\begin{itemize}
    \item explain what a mathematical model is;
    \item distinguish a real-world system from its mathematical model;
    \item identify model objectives;
    \item define state variables, parameters, inputs, and outputs;
    \item distinguish endogenous and exogenous quantities;
    \item identify assumptions explicitly;
    \item distinguish mechanistic and empirical models;
    \item distinguish deterministic and stochastic models;
    \item distinguish discrete and continuous models;
    \item distinguish static and dynamic models;
    \item distinguish linear and nonlinear models;
    \item construct basic algebraic models;
    \item construct basic differential-equation models;
    \item construct basic difference-equation models;
    \item understand conservation laws;
    \item formulate balance equations;
    \item understand compartment models;
    \item use dimensional analysis;
    \item check dimensional consistency;
    \item understand nondimensionalization;
    \item recognize scaling laws;
    \item understand parameter estimation;
    \item distinguish calibration and validation;
    \item understand training and testing conceptually;
    \item recognize identifiability issues;
    \item distinguish structural and practical identifiability;
    \item understand sensitivity analysis;
    \item recognize local and global sensitivity;
    \item understand uncertainty quantification conceptually;
    \item distinguish parameter uncertainty and model uncertainty;
    \item understand residuals;
    \item assess model fit;
    \item recognize overfitting and underfitting;
    \item understand parsimony;
    \item recognize Occam-style model selection principles;
    \item understand interpolation and extrapolation;
    \item recognize extrapolation risk;
    \item understand model verification;
    \item understand model validation;
    \item distinguish verification from validation;
    \item understand scenario analysis;
    \item understand optimization within models;
    \item recognize simulation models;
    \item understand Monte Carlo modeling conceptually;
    \item understand equilibrium models;
    \item recognize dynamical models;
    \item understand feedback;
    \item recognize positive and negative feedback;
    \item understand stability conceptually;
    \item recognize thresholds and tipping points;
    \item understand delay effects;
    \item recognize spatial modeling;
    \item understand network-based modeling;
    \item recognize agent-based modeling conceptually;
    \item understand multiscale modeling conceptually;
    \item recognize data-driven modeling;
    \item understand hybrid mechanistic-data-driven models;
    \item construct a modeling workflow;
    \item understand iterative model refinement;
    \item communicate modeling assumptions and limitations;
    \item recognize ethical and practical limitations of models;
    \item connect mathematical modeling with biology, physics, mechanics,
          economics, finance, networks, optimization, probability,
          statistics, and computation.
\end{itemize}
}

%%%%%%%%%%%%%%%%%%%%%%%%%%%%%%%%%%%%%%%%%%%%%%%%%%%%%%%%%%%%
\subsection{What Is a Mathematical Model?}
\label{subsec:what-is-mathematical-model}
%%%%%%%%%%%%%%%%%%%%%%%%%%%%%%%%%%%%%%%%%%%%%%%%%%%%%%%%%%%%

A mathematical model is a mathematical representation of selected features
of a real or conceptual system.

A model may use:

\[
\boxed{
\text{equations}
}
\]

\[
\boxed{
\text{functions}
}
\]

\[
\boxed{
\text{probability distributions}
}
\]

\[
\boxed{
\text{graphs and networks}
}
\]

or

\[
\boxed{
\text{algorithms}
}
\]

to describe relationships among quantities of interest.

A model is not the system itself.

It is an abstraction designed for a specific purpose.

%%%%%%%%%%%%%%%%%%%%%%%%%%%%%%%%%%%%%%%%%%%%%%%%%%%%%%%%%%%%
\subsection{System Versus Model}
\label{subsec:system-versus-model}
%%%%%%%%%%%%%%%%%%%%%%%%%%%%%%%%%%%%%%%%%%%%%%%%%%%%%%%%%%%%

Suppose a population grows over time.

The real population contains enormous detail:

\[
\boxed{
\text{individual organisms}
}
\]

\[
\boxed{
\text{ages}
}
\]

\[
\boxed{
\text{locations}
}
\]

\[
\boxed{
\text{resources}
}
\]

\[
\boxed{
\text{disease}
}
\]

and many other factors.

A simple model might keep only total population size

\[
\boxed{
P(t).
}
\]

The model intentionally ignores details that are not needed for the current
question.

%%%%%%%%%%%%%%%%%%%%%%%%%%%%%%%%%%%%%%%%%%%%%%%%%%%%%%%%%%%%
\subsection{Purpose Determines the Model}
\label{subsec:model-purpose}
%%%%%%%%%%%%%%%%%%%%%%%%%%%%%%%%%%%%%%%%%%%%%%%%%%%%%%%%%%%%

There is rarely one universally correct model.

The appropriate model depends on the objective.

For example, one population model may be designed to predict:

\[
\boxed{
\text{short-term population size},
}
\]

while another may study:

\[
\boxed{
\text{long-term stability},
}
\]

and another may estimate:

\[
\boxed{
\text{disease transmission risk}.
}
\]

Different questions require different variables and assumptions.

%%%%%%%%%%%%%%%%%%%%%%%%%%%%%%%%%%%%%%%%%%%%%%%%%%%%%%%%%%%%
\subsection{Modeling as Abstraction}
\label{subsec:modeling-abstraction}
%%%%%%%%%%%%%%%%%%%%%%%%%%%%%%%%%%%%%%%%%%%%%%%%%%%%%%%%%%%%

Modeling requires choosing what to include and what to ignore.

The process is:

\[
\boxed{
\text{real system}
\longrightarrow
\text{selected features}
\longrightarrow
\text{mathematical structure}.
}
\]

A good model retains the features needed for the question while avoiding
unnecessary complexity.

%%%%%%%%%%%%%%%%%%%%%%%%%%%%%%%%%%%%%%%%%%%%%%%%%%%%%%%%%%%%
\subsection{Variables}
\label{subsec:modeling-variables}
%%%%%%%%%%%%%%%%%%%%%%%%%%%%%%%%%%%%%%%%%%%%%%%%%%%%%%%%%%%%

A variable represents a quantity that may change.

Examples include:

\[
\boxed{
P(t)
=
\text{population at time }t,
}
\]

\[
\boxed{
T(x,t)
=
\text{temperature at position }x\text{ and time }t,
}
\]

and

\[
\boxed{
S(t)
=
\text{stock price at time }t.
}
\]

Variables are often the primary unknowns of a model.

%%%%%%%%%%%%%%%%%%%%%%%%%%%%%%%%%%%%%%%%%%%%%%%%%%%%%%%%%%%%
\subsection{State Variables}
\label{subsec:state-variables}
%%%%%%%%%%%%%%%%%%%%%%%%%%%%%%%%%%%%%%%%%%%%%%%%%%%%%%%%%%%%

A state variable describes the current condition of a system.

A state vector may be written

\[
\boxed{
\mathbf{x}(t)
=
\begin{bmatrix}
x_1(t)\\
\vdots\\
x_n(t)
\end{bmatrix}.
}
\]

If the model is Markovian, the present state contains enough information to
determine future evolution once inputs and parameters are specified.

%%%%%%%%%%%%%%%%%%%%%%%%%%%%%%%%%%%%%%%%%%%%%%%%%%%%%%%%%%%%
\subsection{Parameters}
\label{subsec:modeling-parameters}
%%%%%%%%%%%%%%%%%%%%%%%%%%%%%%%%%%%%%%%%%%%%%%%%%%%%%%%%%%%%

Parameters are quantities treated as fixed during a particular model run.

For example,

\[
\boxed{
r
=
\text{growth rate},
}
\]

\[
\boxed{
K
=
\text{carrying capacity},
}
\]

and

\[
\boxed{
\beta
=
\text{transmission parameter}.
}
\]

Parameters may be known, estimated from data, or varied during sensitivity
analysis.

%%%%%%%%%%%%%%%%%%%%%%%%%%%%%%%%%%%%%%%%%%%%%%%%%%%%%%%%%%%%
\subsection{Inputs and Outputs}
\label{subsec:model-input-output}
%%%%%%%%%%%%%%%%%%%%%%%%%%%%%%%%%%%%%%%%%%%%%%%%%%%%%%%%%%%%

An input is an externally specified quantity entering the model.

An output is a quantity produced by the model.

A general representation is

\[
\boxed{
\mathbf{y}
=
F(\mathbf{x},\boldsymbol{\theta},\mathbf{u}),
}
\]

where:

\[
\mathbf{x}
=
\text{state or explanatory variables},
\]

\[
\boldsymbol{\theta}
=
\text{parameters},
\]

and

\[
\mathbf{u}
=
\text{external inputs}.
\]

%%%%%%%%%%%%%%%%%%%%%%%%%%%%%%%%%%%%%%%%%%%%%%%%%%%%%%%%%%%%
\subsection{Endogenous and Exogenous Quantities}
\label{subsec:endogenous-exogenous}
%%%%%%%%%%%%%%%%%%%%%%%%%%%%%%%%%%%%%%%%%%%%%%%%%%%%%%%%%%%%

An endogenous variable is determined within the model.

An exogenous variable is imposed from outside the model.

For example, in an economic demand model,

\[
Q=D(P,I),
\]

quantity demanded \(Q\) may be endogenous, while income \(I\) may be treated
as externally specified.

The distinction depends on the modeling boundary.

%%%%%%%%%%%%%%%%%%%%%%%%%%%%%%%%%%%%%%%%%%%%%%%%%%%%%%%%%%%%
\subsection{Assumptions}
\label{subsec:modeling-assumptions}
%%%%%%%%%%%%%%%%%%%%%%%%%%%%%%%%%%%%%%%%%%%%%%%%%%%%%%%%%%%%

Every mathematical model relies on assumptions.

Examples include:

\[
\boxed{
\text{constant parameters},
}
\]

\[
\boxed{
\text{homogeneous populations},
}
\]

\[
\boxed{
\text{perfect mixing},
}
\]

\[
\boxed{
\text{negligible friction},
}
\]

or

\[
\boxed{
\text{independent observations}.
}
\]

Assumptions should be stated explicitly because they determine when the
model is meaningful.

%%%%%%%%%%%%%%%%%%%%%%%%%%%%%%%%%%%%%%%%%%%%%%%%%%%%%%%%%%%%
\subsection{Mechanistic Models}
\label{subsec:mechanistic-models}
%%%%%%%%%%%%%%%%%%%%%%%%%%%%%%%%%%%%%%%%%%%%%%%%%%%%%%%%%%%%

A mechanistic model is built from hypothesized mechanisms governing the
system.

For example,

\[
\boxed{
F=ma
}
\]

relates force, mass, and acceleration through a physical law.

Likewise, an epidemic model may represent infection through contact
mechanisms.

Mechanistic models aim to describe how the system works.

%%%%%%%%%%%%%%%%%%%%%%%%%%%%%%%%%%%%%%%%%%%%%%%%%%%%%%%%%%%%
\subsection{Empirical Models}
\label{subsec:empirical-models}
%%%%%%%%%%%%%%%%%%%%%%%%%%%%%%%%%%%%%%%%%%%%%%%%%%%%%%%%%%%%

An empirical model is built primarily from observed data relationships.

For example,

\[
\boxed{
y
=
\beta_0+\beta_1x+\varepsilon
}
\]

may describe an observed linear association without asserting that the
relationship is a fundamental mechanism.

Empirical models emphasize predictive or descriptive fit.

%%%%%%%%%%%%%%%%%%%%%%%%%%%%%%%%%%%%%%%%%%%%%%%%%%%%%%%%%%%%
\subsection{Hybrid Models}
\label{subsec:hybrid-models}
%%%%%%%%%%%%%%%%%%%%%%%%%%%%%%%%%%%%%%%%%%%%%%%%%%%%%%%%%%%%

Many modern models combine mechanistic equations with data-driven
components.

A system may take form

\[
\boxed{
\frac{d\mathbf{x}}{dt}
=
F(\mathbf{x},\theta)
+
G_{\text{data}}(\mathbf{x}),
}
\]

where \(F\) represents known mechanisms and \(G_{\text{data}}\) represents
an empirically learned correction.

Such hybrid models combine prior scientific structure with observed data.

%%%%%%%%%%%%%%%%%%%%%%%%%%%%%%%%%%%%%%%%%%%%%%%%%%%%%%%%%%%%
\subsection{Deterministic Models}
\label{subsec:deterministic-models}
%%%%%%%%%%%%%%%%%%%%%%%%%%%%%%%%%%%%%%%%%%%%%%%%%%%%%%%%%%%%

A deterministic model produces the same output whenever the same inputs,
parameters, and initial conditions are used.

For example,

\[
\boxed{
\frac{dP}{dt}=rP
}
\]

with fixed \(P(0)\) and \(r\) has a uniquely determined solution under
ordinary conditions.

%%%%%%%%%%%%%%%%%%%%%%%%%%%%%%%%%%%%%%%%%%%%%%%%%%%%%%%%%%%%
\subsection{Stochastic Models}
\label{subsec:stochastic-models}
%%%%%%%%%%%%%%%%%%%%%%%%%%%%%%%%%%%%%%%%%%%%%%%%%%%%%%%%%%%%

A stochastic model includes random quantities.

For example,

\[
\boxed{
X_{t+1}
=
aX_t+\varepsilon_t,
}
\]

where

\[
\varepsilon_t
\]

is random noise.

The output is then described through probability distributions rather than
one deterministic trajectory.

%%%%%%%%%%%%%%%%%%%%%%%%%%%%%%%%%%%%%%%%%%%%%%%%%%%%%%%%%%%%
\subsection{Discrete and Continuous Models}
\label{subsec:discrete-continuous-models}
%%%%%%%%%%%%%%%%%%%%%%%%%%%%%%%%%%%%%%%%%%%%%%%%%%%%%%%%%%%%

A discrete-time model evolves at times

\[
\boxed{
t=0,1,2,\ldots.
}
\]

For example,

\[
\boxed{
x_{n+1}=F(x_n).
}
\]

A continuous-time model evolves over a continuum of times:

\[
\boxed{
\frac{dx}{dt}=F(x,t).
}
\]

The choice depends on the nature and resolution of the system.

%%%%%%%%%%%%%%%%%%%%%%%%%%%%%%%%%%%%%%%%%%%%%%%%%%%%%%%%%%%%
\subsection{Static and Dynamic Models}
\label{subsec:static-dynamic-models}
%%%%%%%%%%%%%%%%%%%%%%%%%%%%%%%%%%%%%%%%%%%%%%%%%%%%%%%%%%%%

A static model does not explicitly describe temporal evolution.

For example,

\[
\boxed{
Q_d=a-bP
}
\]

may relate demand to price at one state.

A dynamic model describes change over time:

\[
\boxed{
\frac{dx}{dt}=F(x,t).
}
\]

%%%%%%%%%%%%%%%%%%%%%%%%%%%%%%%%%%%%%%%%%%%%%%%%%%%%%%%%%%%%
\subsection{Linear and Nonlinear Models}
\label{subsec:linear-nonlinear-models}
%%%%%%%%%%%%%%%%%%%%%%%%%%%%%%%%%%%%%%%%%%%%%%%%%%%%%%%%%%%%

A linear model obeys linear relationships in its relevant unknowns.

For example,

\[
\boxed{
A\mathbf{x}=\mathbf{b}.
}
\]

A nonlinear model may contain:

\[
\boxed{
x^2,
}
\]

\[
\boxed{
xy,
}
\]

\[
\boxed{
e^x,
}
\]

or other nonlinear expressions.

Nonlinear models can exhibit behaviors impossible in linear systems,
including multiple equilibria, bifurcations, and chaos.

%%%%%%%%%%%%%%%%%%%%%%%%%%%%%%%%%%%%%%%%%%%%%%%%%%%%%%%%%%%%
\subsection{A Basic Modeling Cycle}
\label{subsec:modeling-cycle}
%%%%%%%%%%%%%%%%%%%%%%%%%%%%%%%%%%%%%%%%%%%%%%%%%%%%%%%%%%%%

A typical modeling process is iterative:

\[
\boxed{
\text{question}
}
\]

\[
\downarrow
\]

\[
\boxed{
\text{assumptions and variables}
}
\]

\[
\downarrow
\]

\[
\boxed{
\text{mathematical formulation}
}
\]

\[
\downarrow
\]

\[
\boxed{
\text{analysis or simulation}
}
\]

\[
\downarrow
\]

\[
\boxed{
\text{comparison with data}
}
\]

\[
\downarrow
\]

\[
\boxed{
\text{revision}.
}
\]

Models are usually refined rather than created perfectly in one step.

%%%%%%%%%%%%%%%%%%%%%%%%%%%%%%%%%%%%%%%%%%%%%%%%%%%%%%%%%%%%
\subsection{Step 1: Define the Question}
\label{subsec:modeling-define-question}
%%%%%%%%%%%%%%%%%%%%%%%%%%%%%%%%%%%%%%%%%%%%%%%%%%%%%%%%%%%%

A model should begin with a clear objective.

Examples include:

\[
\boxed{
\text{predict future values},
}
\]

\[
\boxed{
\text{explain observed behavior},
}
\]

\[
\boxed{
\text{estimate unknown parameters},
}
\]

\[
\boxed{
\text{compare interventions},
}
\]

or

\[
\boxed{
\text{optimize a decision}.
}
\]

Without a clear question, it is difficult to judge whether a model is useful.

%%%%%%%%%%%%%%%%%%%%%%%%%%%%%%%%%%%%%%%%%%%%%%%%%%%%%%%%%%%%
\subsection{Step 2: Define the System Boundary}
\label{subsec:model-boundary}
%%%%%%%%%%%%%%%%%%%%%%%%%%%%%%%%%%%%%%%%%%%%%%%%%%%%%%%%%%%%

The system boundary determines what is modeled internally and what is
treated as external.

For example, in a lake pollution model one may include:

\[
\boxed{
\text{pollutant concentration}
}
\]

and

\[
\boxed{
\text{water inflow and outflow},
}
\]

while treating weather as an external input.

Changing the boundary changes the model.

%%%%%%%%%%%%%%%%%%%%%%%%%%%%%%%%%%%%%%%%%%%%%%%%%%%%%%%%%%%%
\subsection{Step 3: Choose Variables}
\label{subsec:model-variable-selection}
%%%%%%%%%%%%%%%%%%%%%%%%%%%%%%%%%%%%%%%%%%%%%%%%%%%%%%%%%%%%

Variables should be chosen to capture the information necessary for the
model objective.

Too few variables may omit essential behavior.

Too many variables may produce unnecessary complexity and unidentifiable
parameters.

The goal is not maximum detail.

The goal is sufficient detail.

%%%%%%%%%%%%%%%%%%%%%%%%%%%%%%%%%%%%%%%%%%%%%%%%%%%%%%%%%%%%
\subsection{Step 4: State Assumptions}
\label{subsec:model-assumption-selection}
%%%%%%%%%%%%%%%%%%%%%%%%%%%%%%%%%%%%%%%%%%%%%%%%%%%%%%%%%%%%

Assumptions simplify reality.

For example, one might assume:

\[
\boxed{
\text{constant temperature},
}
\]

\[
\boxed{
\text{constant population},
}
\]

or

\[
\boxed{
\text{uniform spatial mixing}.
}
\]

Each assumption should have a reason and a range of expected validity.

%%%%%%%%%%%%%%%%%%%%%%%%%%%%%%%%%%%%%%%%%%%%%%%%%%%%%%%%%%%%
\subsection{Step 5: Formulate the Mathematics}
\label{subsec:model-formulate-math}
%%%%%%%%%%%%%%%%%%%%%%%%%%%%%%%%%%%%%%%%%%%%%%%%%%%%%%%%%%%%

Translate system relationships into mathematical expressions.

Possible forms include:

\[
\boxed{
\text{algebraic equations},
}
\]

\[
\boxed{
\text{difference equations},
}
\]

\[
\boxed{
\text{differential equations},
}
\]

\[
\boxed{
\text{probability models},
}
\]

\[
\boxed{
\text{optimization problems},
}
\]

or combinations of these.

%%%%%%%%%%%%%%%%%%%%%%%%%%%%%%%%%%%%%%%%%%%%%%%%%%%%%%%%%%%%
\subsection{Step 6: Analyze or Simulate}
\label{subsec:model-analysis-simulation}
%%%%%%%%%%%%%%%%%%%%%%%%%%%%%%%%%%%%%%%%%%%%%%%%%%%%%%%%%%%%

Once formulated, a model may be studied analytically or numerically.

Analytical methods may produce:

\[
\boxed{
\text{closed-form solutions},
}
\]

\[
\boxed{
\text{equilibria},
}
\]

or

\[
\boxed{
\text{stability conditions}.
}
\]

Numerical methods may approximate trajectories or distributions.

%%%%%%%%%%%%%%%%%%%%%%%%%%%%%%%%%%%%%%%%%%%%%%%%%%%%%%%%%%%%
\subsection{Step 7: Compare With Reality}
\label{subsec:model-compare-reality}
%%%%%%%%%%%%%%%%%%%%%%%%%%%%%%%%%%%%%%%%%%%%%%%%%%%%%%%%%%%%

Model predictions should be compared with observations when data are
available.

Possible measures include:

\[
\boxed{
\text{residual error},
}
\]

\[
\boxed{
\text{prediction error},
}
\]

or

\[
\boxed{
\text{likelihood}.
}
\]

Disagreement may reveal poor parameters, poor assumptions, or missing
mechanisms.

%%%%%%%%%%%%%%%%%%%%%%%%%%%%%%%%%%%%%%%%%%%%%%%%%%%%%%%%%%%%
\subsection{Step 8: Revise}
\label{subsec:model-revision}
%%%%%%%%%%%%%%%%%%%%%%%%%%%%%%%%%%%%%%%%%%%%%%%%%%%%%%%%%%%%

A model should be revised when it fails to meet its intended purpose.

Possible revisions include:

\[
\boxed{
\text{changing parameters},
}
\]

\[
\boxed{
\text{adding variables},
}
\]

\[
\boxed{
\text{removing unnecessary detail},
}
\]

or

\[
\boxed{
\text{changing model structure}.
}
\]

Modeling is fundamentally iterative.

%%%%%%%%%%%%%%%%%%%%%%%%%%%%%%%%%%%%%%%%%%%%%%%%%%%%%%%%%%%%
\subsection{Algebraic Models}
\label{subsec:algebraic-models}
%%%%%%%%%%%%%%%%%%%%%%%%%%%%%%%%%%%%%%%%%%%%%%%%%%%%%%%%%%%%

Some models require no explicit time dynamics.

For example, total cost may be modeled as

\[
\boxed{
C(q)=F+cq,
}
\]

where:

\[
F
=
\text{fixed cost},
\]

\[
c
=
\text{unit variable cost},
\]

and

\[
q
=
\text{production quantity}.
\]

This simple model can support optimization and break-even analysis.

%%%%%%%%%%%%%%%%%%%%%%%%%%%%%%%%%%%%%%%%%%%%%%%%%%%%%%%%%%%%
\subsection{Worked Example: Break-Even Model}
\label{subsec:worked-break-even-model}
%%%%%%%%%%%%%%%%%%%%%%%%%%%%%%%%%%%%%%%%%%%%%%%%%%%%%%%%%%%%

\begin{workedexamplebox}{Finding the Break-Even Quantity}

Suppose cost is

\[
C(q)=1000+8q
\]

and revenue is

\[
R(q)=13q.
\]

Break-even occurs when

\[
R(q)=C(q).
\]

Therefore,

\[
13q
=
1000+8q.
\]

Hence

\[
5q=1000,
\]

so

\[
q=200.
\]

\begin{answerbox}
\[
\boxed{
q=200
}
\]

units is the break-even production level.
\end{answerbox}

\end{workedexamplebox}

%%%%%%%%%%%%%%%%%%%%%%%%%%%%%%%%%%%%%%%%%%%%%%%%%%%%%%%%%%%%
\subsection{Difference-Equation Models}
\label{subsec:difference-equation-models}
%%%%%%%%%%%%%%%%%%%%%%%%%%%%%%%%%%%%%%%%%%%%%%%%%%%%%%%%%%%%

A discrete dynamic model may have form

\[
\boxed{
x_{n+1}=F(x_n).
}
\]

For example, constant proportional growth gives

\[
\boxed{
P_{n+1}
=
(1+r)P_n.
}
\]

Iterating gives

\[
\boxed{
P_n
=
P_0(1+r)^n.
}
\]

%%%%%%%%%%%%%%%%%%%%%%%%%%%%%%%%%%%%%%%%%%%%%%%%%%%%%%%%%%%%
\subsection{Differential-Equation Models}
\label{subsec:differential-equation-models}
%%%%%%%%%%%%%%%%%%%%%%%%%%%%%%%%%%%%%%%%%%%%%%%%%%%%%%%%%%%%

Continuous dynamic models often describe rates of change.

A general system is

\[
\boxed{
\frac{d\mathbf{x}}{dt}
=
F(\mathbf{x},t,\theta).
}
\]

The derivative represents how quickly the state changes.

Differential equations are central in biology, physics, engineering, and
economics.

%%%%%%%%%%%%%%%%%%%%%%%%%%%%%%%%%%%%%%%%%%%%%%%%%%%%%%%%%%%%
\subsection{Exponential Growth Model}
\label{subsec:exponential-growth-modeling}
%%%%%%%%%%%%%%%%%%%%%%%%%%%%%%%%%%%%%%%%%%%%%%%%%%%%%%%%%%%%

Assume population growth rate is proportional to current population:

\[
\boxed{
\frac{dP}{dt}=rP.
}
\]

Solving gives

\[
\boxed{
P(t)=P_0e^{rt}.
}
\]

For \(r>0\), population grows exponentially.

For \(r<0\), it decays exponentially.

%%%%%%%%%%%%%%%%%%%%%%%%%%%%%%%%%%%%%%%%%%%%%%%%%%%%%%%%%%%%
\subsection{Worked Example: Exponential Growth}
\label{subsec:worked-exponential-growth-model}
%%%%%%%%%%%%%%%%%%%%%%%%%%%%%%%%%%%%%%%%%%%%%%%%%%%%%%%%%%%%

\begin{workedexamplebox}{Population Projection}

Suppose

\[
P(0)=500
\]

and

\[
r=0.04.
\]

Then

\[
P(t)=500e^{0.04t}.
\]

At \(t=10\),

\[
P(10)
=
500e^{0.4}.
\]

Numerically,

\[
P(10)
\approx
745.9.
\]

\begin{answerbox}
\[
\boxed{
P(10)\approx746.
}
\]
\end{answerbox}

\end{workedexamplebox}

%%%%%%%%%%%%%%%%%%%%%%%%%%%%%%%%%%%%%%%%%%%%%%%%%%%%%%%%%%%%
\subsection{Logistic Growth}
\label{subsec:logistic-growth-modeling}
%%%%%%%%%%%%%%%%%%%%%%%%%%%%%%%%%%%%%%%%%%%%%%%%%%%%%%%%%%%%

Exponential growth cannot continue indefinitely under finite resources.

A logistic model introduces carrying capacity \(K\):

\[
\boxed{
\frac{dP}{dt}
=
rP
\left(
1-\frac{P}{K}
\right).
}
\]

When \(P\ll K\), growth is approximately exponential.

As \(P\) approaches \(K\), growth slows.

%%%%%%%%%%%%%%%%%%%%%%%%%%%%%%%%%%%%%%%%%%%%%%%%%%%%%%%%%%%%
\subsection{Equilibria of the Logistic Model}
\label{subsec:logistic-equilibria}
%%%%%%%%%%%%%%%%%%%%%%%%%%%%%%%%%%%%%%%%%%%%%%%%%%%%%%%%%%%%

Equilibria satisfy

\[
0
=
rP
\left(
1-\frac{P}{K}
\right).
\]

Therefore,

\[
\boxed{
P^*=0
}
\]

or

\[
\boxed{
P^*=K.
}
\]

For \(r>0\), the positive equilibrium \(P=K\) is stable under the standard
model.

%%%%%%%%%%%%%%%%%%%%%%%%%%%%%%%%%%%%%%%%%%%%%%%%%%%%%%%%%%%%
\subsection{Worked Example: Logistic Rate}
\label{subsec:worked-logistic-rate}
%%%%%%%%%%%%%%%%%%%%%%%%%%%%%%%%%%%%%%%%%%%%%%%%%%%%%%%%%%%%

\begin{workedexamplebox}{Growth at Half Carrying Capacity}

For

\[
\frac{dP}{dt}
=
rP
\left(
1-\frac{P}{K}
\right),
\]

let

\[
P=\frac{K}{2}.
\]

Then

\[
\frac{dP}{dt}
=
r\frac{K}{2}
\left(
1-\frac12
\right).
\]

Thus

\[
\begin{answerbox}
\[
\boxed{
\frac{dP}{dt}
=
\frac{rK}{4}.
}
\]
\end{answerbox}

This is the maximum growth rate of the standard logistic growth term.

\end{workedexamplebox}

%%%%%%%%%%%%%%%%%%%%%%%%%%%%%%%%%%%%%%%%%%%%%%%%%%%%%%%%%%%%
\subsection{Conservation Laws}
\label{subsec:modeling-conservation-laws}
%%%%%%%%%%%%%%%%%%%%%%%%%%%%%%%%%%%%%%%%%%%%%%%%%%%%%%%%%%%%

Many models are built from conservation principles.

A generic balance law is

\[
\boxed{
\text{accumulation}
=
\text{input}
-
\text{output}
+
\text{generation}
-
\text{consumption}.
}
\]

This principle appears in models of:

\[
\boxed{
\text{mass},
}
\]

\[
\boxed{
\text{energy},
}
\]

\[
\boxed{
\text{population},
}
\]

and

\[
\boxed{
\text{financial accounts}.
}
\]

%%%%%%%%%%%%%%%%%%%%%%%%%%%%%%%%%%%%%%%%%%%%%%%%%%%%%%%%%%%%
\subsection{Balance Equations}
\label{subsec:balance-equations}
%%%%%%%%%%%%%%%%%%%%%%%%%%%%%%%%%%%%%%%%%%%%%%%%%%%%%%%%%%%%

Let \(Q(t)\) denote the amount of some conserved quantity.

Then a general balance equation is

\[
\boxed{
\frac{dQ}{dt}
=
R_{\text{in}}
-
R_{\text{out}}
+
R_{\text{source}}
-
R_{\text{sink}}.
}
\]

This provides a systematic method for building many differential-equation
models.

%%%%%%%%%%%%%%%%%%%%%%%%%%%%%%%%%%%%%%%%%%%%%%%%%%%%%%%%%%%%
\subsection{Worked Example: Tank Mixing}
\label{subsec:worked-tank-mixing}
%%%%%%%%%%%%%%%%%%%%%%%%%%%%%%%%%%%%%%%%%%%%%%%%%%%%%%%%%%%%

\begin{workedexamplebox}{Salt in a Well-Mixed Tank}

Let \(Q(t)\) be the amount of salt in a tank of constant volume \(V\).

Suppose liquid enters at rate \(r\) with salt concentration \(c_{\text{in}}\)
and leaves at the same rate.

Under perfect mixing, outflow concentration is

\[
\frac{Q(t)}{V}.
\]

Thus

\[
\frac{dQ}{dt}
=
rc_{\text{in}}
-
r\frac{Q}{V}.
\]

\begin{answerbox}
\[
\boxed{
\frac{dQ}{dt}
=
r
\left(
c_{\text{in}}
-
\frac{Q}{V}
\right).
}
\]
\end{answerbox}

\end{workedexamplebox}

%%%%%%%%%%%%%%%%%%%%%%%%%%%%%%%%%%%%%%%%%%%%%%%%%%%%%%%%%%%%
\subsection{Compartment Models}
\label{subsec:compartment-models}
%%%%%%%%%%%%%%%%%%%%%%%%%%%%%%%%%%%%%%%%%%%%%%%%%%%%%%%%%%%%

A compartment model divides a system into interacting groups or regions.

If

\[
x_i(t)
\]

is the amount in compartment \(i\), then

\[
\boxed{
\frac{dx_i}{dt}
=
\text{flow into }i
-
\text{flow out of }i.
}
\]

Compartment models are common in:

\[
\boxed{
\text{epidemiology},
}
\]

\[
\boxed{
\text{pharmacokinetics},
}
\]

\[
\boxed{
\text{ecology},
}
\]

and

\[
\boxed{
\text{transport systems}.
}
\]

%%%%%%%%%%%%%%%%%%%%%%%%%%%%%%%%%%%%%%%%%%%%%%%%%%%%%%%%%%%%
\subsection{A Two-Compartment Model}
\label{subsec:two-compartment-model}
%%%%%%%%%%%%%%%%%%%%%%%%%%%%%%%%%%%%%%%%%%%%%%%%%%%%%%%%%%%%

Suppose material moves between two compartments.

Let

\[
x_1(t),x_2(t)
\]

be their amounts.

With transfer rates \(k_{12}\) and \(k_{21}\),

\[
\boxed{
\frac{dx_1}{dt}
=
-k_{12}x_1+k_{21}x_2,
}
\]

\[
\boxed{
\frac{dx_2}{dt}
=
k_{12}x_1-k_{21}x_2.
}
\]

Adding gives

\[
\boxed{
\frac{d}{dt}(x_1+x_2)=0,
}
\]

so total material is conserved.

%%%%%%%%%%%%%%%%%%%%%%%%%%%%%%%%%%%%%%%%%%%%%%%%%%%%%%%%%%%%
\subsection{Dimensional Analysis}
\label{subsec:dimensional-analysis-modeling}
%%%%%%%%%%%%%%%%%%%%%%%%%%%%%%%%%%%%%%%%%%%%%%%%%%%%%%%%%%%%

Dimensional analysis checks whether mathematical expressions are compatible
with physical units.

Suppose

\[
x
\]

has units of length and

\[
t
\]

has units of time.

Then velocity

\[
\frac{dx}{dt}
\]

has units

\[
\boxed{
\frac{\text{length}}{\text{time}}.
}
\]

Every term in a valid physical equation must have compatible dimensions.

%%%%%%%%%%%%%%%%%%%%%%%%%%%%%%%%%%%%%%%%%%%%%%%%%%%%%%%%%%%%
\subsection{Dimensional Consistency}
\label{subsec:dimensional-consistency}
%%%%%%%%%%%%%%%%%%%%%%%%%%%%%%%%%%%%%%%%%%%%%%%%%%%%%%%%%%%%

Consider

\[
x(t)=x_0+vt.
\]

The units are:

\[
[x]=L,
\]

\[
[x_0]=L,
\]

\[
[v]=LT^{-1},
\]

and

\[
[t]=T.
\]

Thus

\[
[vt]
=
LT^{-1}T
=
L.
\]

Both terms on the right have units of length.

The equation is dimensionally consistent.

%%%%%%%%%%%%%%%%%%%%%%%%%%%%%%%%%%%%%%%%%%%%%%%%%%%%%%%%%%%%
\subsection{Worked Example: Detecting a Unit Error}
\label{subsec:worked-dimensional-error}
%%%%%%%%%%%%%%%%%%%%%%%%%%%%%%%%%%%%%%%%%%%%%%%%%%%%%%%%%%%%

\begin{workedexamplebox}{Checking a Motion Formula}

Suppose someone proposes

\[
x=x_0+vt+at.
\]

The units of the three terms are

\[
[x_0]=L,
\]

\[
[vt]=L,
\]

but

\[
[at]
=
LT^{-2}T
=
LT^{-1}.
\]

Thus \(at\) has units of velocity, not position.

\begin{answerbox}
\[
\boxed{
x=x_0+vt+at
\text{ is dimensionally inconsistent}.
}
\]
\end{answerbox}

A correct constant-acceleration position term is proportional to \(at^2\).

\end{workedexamplebox}

%%%%%%%%%%%%%%%%%%%%%%%%%%%%%%%%%%%%%%%%%%%%%%%%%%%%%%%%%%%%
\subsection{Dimensionless Quantities}
\label{subsec:dimensionless-quantities}
%%%%%%%%%%%%%%%%%%%%%%%%%%%%%%%%%%%%%%%%%%%%%%%%%%%%%%%%%%%%

A dimensionless quantity has no physical units.

For example, the ratio

\[
\boxed{
\frac{x}{L}
}
\]

is dimensionless when both \(x\) and \(L\) are lengths.

Dimensionless groups often reveal the true parameters controlling system
behavior.

%%%%%%%%%%%%%%%%%%%%%%%%%%%%%%%%%%%%%%%%%%%%%%%%%%%%%%%%%%%%
\subsection{Nondimensionalization}
\label{subsec:nondimensionalization}
%%%%%%%%%%%%%%%%%%%%%%%%%%%%%%%%%%%%%%%%%%%%%%%%%%%%%%%%%%%%

Nondimensionalization rescales variables.

Suppose characteristic scales are:

\[
\boxed{
x=L\hat{x},
}
\]

and

\[
\boxed{
t=T\hat{t}.
}
\]

Then the dimensionless variables are

\[
\boxed{
\hat{x}=\frac{x}{L},
\qquad
\hat{t}=\frac{t}{T}.
}
\]

Substituting these into the model can reduce the number of independent
parameters.

%%%%%%%%%%%%%%%%%%%%%%%%%%%%%%%%%%%%%%%%%%%%%%%%%%%%%%%%%%%%
\subsection{Worked Example: Nondimensional Logistic Equation}
\label{subsec:worked-nondimensional-logistic}
%%%%%%%%%%%%%%%%%%%%%%%%%%%%%%%%%%%%%%%%%%%%%%%%%%%%%%%%%%%%

\begin{workedexamplebox}{Scaling Population and Time}

Start with

\[
\frac{dP}{dt}
=
rP
\left(
1-\frac{P}{K}
\right).
\]

Define

\[
u=\frac{P}{K}
\]

and

\[
\tau=rt.
\]

Since

\[
P=Ku,
\]

we have

\[
\frac{dP}{dt}
=
Kr\frac{du}{d\tau}.
\]

Substitute:

\[
Kr\frac{du}{d\tau}
=
rKu(1-u).
\]

Cancel \(rK\):

\[
\begin{answerbox}
\[
\boxed{
\frac{du}{d\tau}
=
u(1-u).
}
\]
\end{answerbox}

The dimensionless model contains no free parameters.

\end{workedexamplebox}

%%%%%%%%%%%%%%%%%%%%%%%%%%%%%%%%%%%%%%%%%%%%%%%%%%%%%%%%%%%%
\subsection{Scaling Laws}
\label{subsec:modeling-scaling-laws}
%%%%%%%%%%%%%%%%%%%%%%%%%%%%%%%%%%%%%%%%%%%%%%%%%%%%%%%%%%%%

A scaling law relates quantities through powers.

For example,

\[
\boxed{
Y
=
cX^\alpha.
}
\]

Taking logarithms gives

\[
\boxed{
\log Y
=
\log c
+
\alpha\log X.
}
\]

Thus power-law relationships become linear on log-log coordinates.

%%%%%%%%%%%%%%%%%%%%%%%%%%%%%%%%%%%%%%%%%%%%%%%%%%%%%%%%%%%%
\subsection{Parameter Estimation}
\label{subsec:model-parameter-estimation}
%%%%%%%%%%%%%%%%%%%%%%%%%%%%%%%%%%%%%%%%%%%%%%%%%%%%%%%%%%%%

Parameters are often unknown and must be estimated from data.

Suppose observations are

\[
\boxed{
y_i
}
\]

and model predictions are

\[
\boxed{
\hat{y}_i(\theta).
}
\]

One may estimate \(\theta\) by minimizing

\[
\boxed{
S(\theta)
=
\sum_{i=1}^{n}
\left(
y_i-\hat{y}_i(\theta)
\right)^2.
}
\]

This is nonlinear least squares in the general case.

%%%%%%%%%%%%%%%%%%%%%%%%%%%%%%%%%%%%%%%%%%%%%%%%%%%%%%%%%%%%
\subsection{Residuals}
\label{subsec:model-residuals}
%%%%%%%%%%%%%%%%%%%%%%%%%%%%%%%%%%%%%%%%%%%%%%%%%%%%%%%%%%%%

A residual is

\[
\boxed{
r_i
=
y_i-\hat{y}_i.
}
\]

Residuals measure differences between observations and model predictions.

A good model should not merely produce small residuals.

Residual patterns should also be examined for systematic structure.

%%%%%%%%%%%%%%%%%%%%%%%%%%%%%%%%%%%%%%%%%%%%%%%%%%%%%%%%%%%%
\subsection{Worked Example: Least-Squares Parameter}
\label{subsec:worked-model-least-squares}
%%%%%%%%%%%%%%%%%%%%%%%%%%%%%%%%%%%%%%%%%%%%%%%%%%%%%%%%%%%%

\begin{workedexamplebox}{Fitting a Proportional Model}

Suppose the model is

\[
y=ax
\]

and data are

\[
(x_1,y_1),\ldots,(x_n,y_n).
\]

Minimize

\[
S(a)
=
\sum_{i=1}^{n}
(y_i-ax_i)^2.
\]

Differentiate:

\[
S'(a)
=
-2
\sum_{i=1}^{n}
x_i(y_i-ax_i).
\]

Set equal to zero:

\[
\sum x_i y_i
-
a\sum x_i^2
=
0.
\]

Hence

\begin{answerbox}
\[
\boxed{
\hat{a}
=
\frac{
\sum_{i=1}^{n}x_i y_i
}{
\sum_{i=1}^{n}x_i^2
}.
}
\]
\end{answerbox}

\end{workedexamplebox}

%%%%%%%%%%%%%%%%%%%%%%%%%%%%%%%%%%%%%%%%%%%%%%%%%%%%%%%%%%%%
\subsection{Calibration}
\label{subsec:model-calibration}
%%%%%%%%%%%%%%%%%%%%%%%%%%%%%%%%%%%%%%%%%%%%%%%%%%%%%%%%%%%%

Calibration is the process of selecting parameter values so the model
agrees sufficiently well with observed data.

Calibration may use:

\[
\boxed{
\text{least squares},
}
\]

\[
\boxed{
\text{maximum likelihood},
}
\]

\[
\boxed{
\text{Bayesian inference},
}
\]

or other optimization procedures.

Calibration does not by itself establish that the model is correct.

%%%%%%%%%%%%%%%%%%%%%%%%%%%%%%%%%%%%%%%%%%%%%%%%%%%%%%%%%%%%
\subsection{Validation}
\label{subsec:model-validation}
%%%%%%%%%%%%%%%%%%%%%%%%%%%%%%%%%%%%%%%%%%%%%%%%%%%%%%%%%%%%

Validation evaluates whether the model adequately represents the system for
its intended use.

This may involve comparing model predictions with data not used in
calibration.

A useful distinction is:

\[
\boxed{
\text{calibration}
=
\text{fit the model},
}
\]

\[
\boxed{
\text{validation}
=
\text{test the model}.
}
\]

%%%%%%%%%%%%%%%%%%%%%%%%%%%%%%%%%%%%%%%%%%%%%%%%%%%%%%%%%%%%
\subsection{Verification}
\label{subsec:model-verification}
%%%%%%%%%%%%%%%%%%%%%%%%%%%%%%%%%%%%%%%%%%%%%%%%%%%%%%%%%%%%

Verification asks whether the mathematical or computational model has been
implemented correctly.

The question is:

\[
\boxed{
\text{Did we solve the equations correctly?}
}
\]

Validation asks:

\[
\boxed{
\text{Are these the right equations for the intended purpose?}
}
\]

These questions must not be confused.

%%%%%%%%%%%%%%%%%%%%%%%%%%%%%%%%%%%%%%%%%%%%%%%%%%%%%%%%%%%%
\subsection{Verification Versus Validation}
\label{subsec:verification-versus-validation}
%%%%%%%%%%%%%%%%%%%%%%%%%%%%%%%%%%%%%%%%%%%%%%%%%%%%%%%%%%%%

A common summary is:

\[
\boxed{
\text{Verification: solving the model right.}
}
\]

\[
\boxed{
\text{Validation: solving the right model.}
}
\]

A simulation can be numerically perfect and still represent reality poorly.

%%%%%%%%%%%%%%%%%%%%%%%%%%%%%%%%%%%%%%%%%%%%%%%%%%%%%%%%%%%%
\subsection{Training and Testing}
\label{subsec:model-training-testing}
%%%%%%%%%%%%%%%%%%%%%%%%%%%%%%%%%%%%%%%%%%%%%%%%%%%%%%%%%%%%

For data-driven models, data are often separated into:

\[
\boxed{
\text{training data}
}
\]

and

\[
\boxed{
\text{testing data}.
}
\]

Training data are used to fit parameters.

Testing data assess performance on observations not used during fitting.

This helps identify overfitting.

%%%%%%%%%%%%%%%%%%%%%%%%%%%%%%%%%%%%%%%%%%%%%%%%%%%%%%%%%%%%
\subsection{Overfitting}
\label{subsec:model-overfitting}
%%%%%%%%%%%%%%%%%%%%%%%%%%%%%%%%%%%%%%%%%%%%%%%%%%%%%%%%%%%%

An overfit model matches training data extremely closely but performs poorly
on new data.

This often occurs when the model is too flexible relative to the amount of
information available.

Conceptually,

\[
\boxed{
\text{excellent fit}
\not\Rightarrow
\text{excellent prediction}.
}
\]

%%%%%%%%%%%%%%%%%%%%%%%%%%%%%%%%%%%%%%%%%%%%%%%%%%%%%%%%%%%%
\subsection{Underfitting}
\label{subsec:model-underfitting}
%%%%%%%%%%%%%%%%%%%%%%%%%%%%%%%%%%%%%%%%%%%%%%%%%%%%%%%%%%%%

An underfit model is too simple to capture important structure in the data.

It may have poor performance on both training and test data.

Thus model complexity must balance:

\[
\boxed{
\text{simplicity}
}
\]

and

\[
\boxed{
\text{adequate representation}.
}
\]

%%%%%%%%%%%%%%%%%%%%%%%%%%%%%%%%%%%%%%%%%%%%%%%%%%%%%%%%%%%%
\subsection{Parsimony}
\label{subsec:model-parsimony}
%%%%%%%%%%%%%%%%%%%%%%%%%%%%%%%%%%%%%%%%%%%%%%%%%%%%%%%%%%%%

Parsimony favors simpler models when additional complexity provides little
meaningful improvement.

A simpler model may be:

\[
\boxed{
\text{easier to interpret},
}
\]

\[
\boxed{
\text{easier to estimate},
}
\]

and

\[
\boxed{
\text{more robust}.
}
\]

Simplicity is not automatically better; it must still meet the modeling
objective.

%%%%%%%%%%%%%%%%%%%%%%%%%%%%%%%%%%%%%%%%%%%%%%%%%%%%%%%%%%%%
\subsection{Interpolation}
\label{subsec:model-interpolation}
%%%%%%%%%%%%%%%%%%%%%%%%%%%%%%%%%%%%%%%%%%%%%%%%%%%%%%%%%%%%

Interpolation predicts within the region covered by observed data.

For example, if measurements exist for

\[
0\leq x\leq10,
\]

predicting at

\[
x=6
\]

is interpolation.

This is usually safer than extrapolation.

%%%%%%%%%%%%%%%%%%%%%%%%%%%%%%%%%%%%%%%%%%%%%%%%%%%%%%%%%%%%
\subsection{Extrapolation}
\label{subsec:model-extrapolation}
%%%%%%%%%%%%%%%%%%%%%%%%%%%%%%%%%%%%%%%%%%%%%%%%%%%%%%%%%%%%

Extrapolation predicts outside the observed range.

For the same data,

\[
x=50
\]

would be an extrapolation.

Even a model that fits observed data very well may fail far outside the
calibration region.

%%%%%%%%%%%%%%%%%%%%%%%%%%%%%%%%%%%%%%%%%%%%%%%%%%%%%%%%%%%%
\subsection{Worked Example: Extrapolation Risk}
\label{subsec:worked-extrapolation-risk}
%%%%%%%%%%%%%%%%%%%%%%%%%%%%%%%%%%%%%%%%%%%%%%%%%%%%%%%%%%%%

\begin{workedexamplebox}{Linear Growth Cannot Continue Forever}

Suppose observed production over a short interval is approximated by

\[
P(t)=100+20t.
\]

The model may fit well for

\[
0\leq t\leq5.
\]

Using the same formula for

\[
t=100
\]

predicts

\[
P(100)=2100.
\]

However, the system may encounter resource, capacity, or market constraints
long before then.

\begin{answerbox}
\[
\boxed{
\text{A good local fit does not justify unlimited extrapolation.}
}
\]
\end{answerbox}

\end{workedexamplebox}

%%%%%%%%%%%%%%%%%%%%%%%%%%%%%%%%%%%%%%%%%%%%%%%%%%%%%%%%%%%%
\subsection{Identifiability}
\label{subsec:model-identifiability}
%%%%%%%%%%%%%%%%%%%%%%%%%%%%%%%%%%%%%%%%%%%%%%%%%%%%%%%%%%%%

A parameter is identifiable if different parameter values lead to
distinguishable model behavior under the available observations.

If

\[
\theta_1\neq\theta_2
\]

but

\[
\boxed{
F(t;\theta_1)
=
F(t;\theta_2)
}
\]

for all observable quantities, then the data cannot distinguish the two
parameter values.

%%%%%%%%%%%%%%%%%%%%%%%%%%%%%%%%%%%%%%%%%%%%%%%%%%%%%%%%%%%%
\subsection{Structural Identifiability}
\label{subsec:structural-identifiability}
%%%%%%%%%%%%%%%%%%%%%%%%%%%%%%%%%%%%%%%%%%%%%%%%%%%%%%%%%%%%

Structural identifiability concerns whether parameters could be uniquely
recovered with ideal, noise-free, sufficiently rich data.

It is a property of the model structure.

If a parameter is structurally unidentifiable, collecting more data of the
same kind cannot solve the problem.

%%%%%%%%%%%%%%%%%%%%%%%%%%%%%%%%%%%%%%%%%%%%%%%%%%%%%%%%%%%%
\subsection{Practical Identifiability}
\label{subsec:practical-identifiability}
%%%%%%%%%%%%%%%%%%%%%%%%%%%%%%%%%%%%%%%%%%%%%%%%%%%%%%%%%%%%

A parameter may be structurally identifiable but difficult to estimate in
practice because of:

\[
\boxed{
\text{measurement noise},
}
\]

\[
\boxed{
\text{limited sample size},
}
\]

or

\[
\boxed{
\text{weak sensitivity}.
}
\]

This is practical identifiability.

%%%%%%%%%%%%%%%%%%%%%%%%%%%%%%%%%%%%%%%%%%%%%%%%%%%%%%%%%%%%
\subsection{Worked Example: Nonidentifiable Product}
\label{subsec:worked-nonidentifiable-product}
%%%%%%%%%%%%%%%%%%%%%%%%%%%%%%%%%%%%%%%%%%%%%%%%%%%%%%%%%%%%

\begin{workedexamplebox}{Two Parameters Appearing Only as a Product}

Suppose observations are described by

\[
y(t)
=
ab\,t.
\]

Only the product

\[
ab
\]

appears.

If

\[
ab=6,
\]

then parameter pairs

\[
(a,b)=(1,6),
\]

\[
(2,3),
\]

and

\[
(3,2)
\]

all produce the same model output.

\begin{answerbox}
\[
\boxed{
a
\text{ and }
b
\text{ are not individually identifiable from this model output.}
}
\]
\end{answerbox}

\end{workedexamplebox}

%%%%%%%%%%%%%%%%%%%%%%%%%%%%%%%%%%%%%%%%%%%%%%%%%%%%%%%%%%%%
\subsection{Sensitivity Analysis}
\label{subsec:model-sensitivity-analysis}
%%%%%%%%%%%%%%%%%%%%%%%%%%%%%%%%%%%%%%%%%%%%%%%%%%%%%%%%%%%%

Sensitivity analysis studies how model outputs change when inputs or
parameters change.

Suppose

\[
y=f(\theta).
\]

A local sensitivity is

\[
\boxed{
\frac{\partial y}{\partial\theta}.
}
\]

Large magnitude indicates strong local response of \(y\) to changes in
\(\theta\).

%%%%%%%%%%%%%%%%%%%%%%%%%%%%%%%%%%%%%%%%%%%%%%%%%%%%%%%%%%%%
\subsection{Normalized Sensitivity}
\label{subsec:normalized-sensitivity}
%%%%%%%%%%%%%%%%%%%%%%%%%%%%%%%%%%%%%%%%%%%%%%%%%%%%%%%%%%%%

Sensitivity depends on units.

A dimensionless normalized sensitivity may be defined by

\[
\boxed{
S_\theta
=
\frac{\theta}{y}
\frac{\partial y}{\partial\theta},
}
\]

when \(y\neq0\) and \(\theta\neq0\).

This approximates the relative change in output per relative change in the
parameter.

%%%%%%%%%%%%%%%%%%%%%%%%%%%%%%%%%%%%%%%%%%%%%%%%%%%%%%%%%%%%
\subsection{Worked Example: Exponential Sensitivity}
\label{subsec:worked-exponential-sensitivity}
%%%%%%%%%%%%%%%%%%%%%%%%%%%%%%%%%%%%%%%%%%%%%%%%%%%%%%%%%%%%

\begin{workedexamplebox}{Sensitivity to Growth Rate}

Consider

\[
P(t)=P_0e^{rt}.
\]

Differentiate with respect to \(r\):

\[
\frac{\partial P}{\partial r}
=
tP_0e^{rt}.
\]

Since

\[
P(t)=P_0e^{rt},
\]

we obtain

\[
\begin{answerbox}
\[
\boxed{
\frac{\partial P}{\partial r}
=
tP(t).
}
\]
\end{answerbox}

Thus sensitivity to \(r\) grows with time.

\end{workedexamplebox}

%%%%%%%%%%%%%%%%%%%%%%%%%%%%%%%%%%%%%%%%%%%%%%%%%%%%%%%%%%%%
\subsection{Local Sensitivity}
\label{subsec:local-sensitivity}
%%%%%%%%%%%%%%%%%%%%%%%%%%%%%%%%%%%%%%%%%%%%%%%%%%%%%%%%%%%%

Local sensitivity examines small changes near one chosen parameter set.

For parameter vector

\[
\boldsymbol{\theta},
\]

one may compute a sensitivity matrix

\[
\boxed{
S_{ij}
=
\frac{\partial y_i}{\partial\theta_j}.
}
\]

This is closely related to a Jacobian matrix.

%%%%%%%%%%%%%%%%%%%%%%%%%%%%%%%%%%%%%%%%%%%%%%%%%%%%%%%%%%%%
\subsection{Global Sensitivity}
\label{subsec:global-sensitivity}
%%%%%%%%%%%%%%%%%%%%%%%%%%%%%%%%%%%%%%%%%%%%%%%%%%%%%%%%%%%%

Global sensitivity explores parameter variation over a larger region.

It may investigate:

\[
\boxed{
\text{parameter interactions},
}
\]

\[
\boxed{
\text{nonlinear effects},
}
\]

and

\[
\boxed{
\text{uncertainty distributions}.
}
\]

Monte Carlo sampling is often used for this purpose.

%%%%%%%%%%%%%%%%%%%%%%%%%%%%%%%%%%%%%%%%%%%%%%%%%%%%%%%%%%%%
\subsection{Uncertainty Quantification}
\label{subsec:uncertainty-quantification-modeling}
%%%%%%%%%%%%%%%%%%%%%%%%%%%%%%%%%%%%%%%%%%%%%%%%%%%%%%%%%%%%

Uncertainty quantification studies how uncertain inputs produce uncertain
outputs.

If

\[
\Theta
\]

is uncertain and

\[
Y=f(\Theta),
\]

then uncertainty in \(\Theta\) induces uncertainty in \(Y\).

The result may be summarized through:

\[
\boxed{
\text{intervals},
}
\]

\[
\boxed{
\text{probability distributions},
}
\]

or

\[
\boxed{
\text{credible or confidence regions}.
}
\]

%%%%%%%%%%%%%%%%%%%%%%%%%%%%%%%%%%%%%%%%%%%%%%%%%%%%%%%%%%%%
\subsection{Sources of Model Uncertainty}
\label{subsec:model-uncertainty-sources}
%%%%%%%%%%%%%%%%%%%%%%%%%%%%%%%%%%%%%%%%%%%%%%%%%%%%%%%%%%%%

Model uncertainty may come from:

\[
\boxed{
\text{measurement error},
}
\]

\[
\boxed{
\text{parameter uncertainty},
}
\]

\[
\boxed{
\text{initial-condition uncertainty},
}
\]

\[
\boxed{
\text{random variation},
}
\]

and

\[
\boxed{
\text{structural model error}.
}
\]

These sources should not automatically be treated as equivalent.

%%%%%%%%%%%%%%%%%%%%%%%%%%%%%%%%%%%%%%%%%%%%%%%%%%%%%%%%%%%%
\subsection{Aleatoric and Epistemic Uncertainty}
\label{subsec:model-aleatoric-epistemic}
%%%%%%%%%%%%%%%%%%%%%%%%%%%%%%%%%%%%%%%%%%%%%%%%%%%%%%%%%%%%

Aleatoric uncertainty represents inherent variability.

Epistemic uncertainty represents incomplete knowledge.

Conceptually,

\[
\boxed{
\text{aleatoric}
=
\text{randomness in the system},
}
\]

while

\[
\boxed{
\text{epistemic}
=
\text{uncertainty about the system or parameters}.
}
\]

The distinction can influence how uncertainty is modeled and reduced.

%%%%%%%%%%%%%%%%%%%%%%%%%%%%%%%%%%%%%%%%%%%%%%%%%%%%%%%%%%%%
\subsection{Scenario Analysis}
\label{subsec:model-scenario-analysis}
%%%%%%%%%%%%%%%%%%%%%%%%%%%%%%%%%%%%%%%%%%%%%%%%%%%%%%%%%%%%

Scenario analysis compares model outcomes under different assumptions.

For example:

\[
\boxed{
\text{baseline},
}
\]

\[
\boxed{
\text{high-growth scenario},
}
\]

\[
\boxed{
\text{low-growth scenario},
}
\]

or

\[
\boxed{
\text{intervention scenario}.
}
\]

Scenarios are not necessarily probability forecasts.

They represent structured hypothetical cases.

%%%%%%%%%%%%%%%%%%%%%%%%%%%%%%%%%%%%%%%%%%%%%%%%%%%%%%%%%%%%
\subsection{Simulation}
\label{subsec:mathematical-simulation}
%%%%%%%%%%%%%%%%%%%%%%%%%%%%%%%%%%%%%%%%%%%%%%%%%%%%%%%%%%%%

Simulation numerically reproduces the behavior implied by a mathematical
model.

A simulation may generate:

\[
\boxed{
\text{trajectories},
}
\]

\[
\boxed{
\text{event sequences},
}
\]

\[
\boxed{
\text{spatial patterns},
}
\]

or

\[
\boxed{
\text{probability distributions}.
}
\]

Simulation is especially useful when exact analysis is impossible.

%%%%%%%%%%%%%%%%%%%%%%%%%%%%%%%%%%%%%%%%%%%%%%%%%%%%%%%%%%%%
\subsection{Monte Carlo Modeling}
\label{subsec:modeling-monte-carlo}
%%%%%%%%%%%%%%%%%%%%%%%%%%%%%%%%%%%%%%%%%%%%%%%%%%%%%%%%%%%%

Monte Carlo methods repeatedly sample random inputs or events.

If

\[
Y=g(X),
\]

and \(X\) has a probability distribution, one may simulate

\[
X_1,\ldots,X_N
\]

and calculate

\[
Y_i=g(X_i).
\]

Then

\[
\boxed{
\frac1N
\sum_{i=1}^{N}Y_i
}
\]

can estimate

\[
\mathbb{E}[Y].
\]

%%%%%%%%%%%%%%%%%%%%%%%%%%%%%%%%%%%%%%%%%%%%%%%%%%%%%%%%%%%%
\subsection{Equilibrium Models}
\label{subsec:model-equilibrium}
%%%%%%%%%%%%%%%%%%%%%%%%%%%%%%%%%%%%%%%%%%%%%%%%%%%%%%%%%%%%

An equilibrium is a state that does not change under the model dynamics.

For

\[
\frac{d\mathbf{x}}{dt}
=
F(\mathbf{x}),
\]

an equilibrium satisfies

\[
\boxed{
F(\mathbf{x}^*)=0.
}
\]

Equilibria provide important long-term reference states.

%%%%%%%%%%%%%%%%%%%%%%%%%%%%%%%%%%%%%%%%%%%%%%%%%%%%%%%%%%%%
\subsection{Feedback}
\label{subsec:model-feedback}
%%%%%%%%%%%%%%%%%%%%%%%%%%%%%%%%%%%%%%%%%%%%%%%%%%%%%%%%%%%%

Feedback occurs when a system's current state influences its future rate of
change.

For example,

\[
\boxed{
\frac{dx}{dt}=f(x)
}
\]

already contains state-dependent feedback.

Feedback can stabilize or destabilize a system.

%%%%%%%%%%%%%%%%%%%%%%%%%%%%%%%%%%%%%%%%%%%%%%%%%%%%%%%%%%%%
\subsection{Negative Feedback}
\label{subsec:negative-feedback}
%%%%%%%%%%%%%%%%%%%%%%%%%%%%%%%%%%%%%%%%%%%%%%%%%%%%%%%%%%%%

Negative feedback tends to oppose deviations.

For example,

\[
\boxed{
\frac{dx}{dt}
=
-k(x-x^*),
\qquad
k>0,
}
\]

drives \(x\) toward equilibrium \(x^*\).

The solution is

\[
\boxed{
x(t)-x^*
=
(x_0-x^*)e^{-kt}.
}
\]

%%%%%%%%%%%%%%%%%%%%%%%%%%%%%%%%%%%%%%%%%%%%%%%%%%%%%%%%%%%%
\subsection{Positive Feedback}
\label{subsec:positive-feedback}
%%%%%%%%%%%%%%%%%%%%%%%%%%%%%%%%%%%%%%%%%%%%%%%%%%%%%%%%%%%%

Positive feedback amplifies deviations.

For example,

\[
\boxed{
\frac{dx}{dt}
=
kx,
\qquad
k>0,
}
\]

gives

\[
x(t)=x_0e^{kt}.
\]

Small positive values grow exponentially.

Positive feedback can create rapid growth or instability.

%%%%%%%%%%%%%%%%%%%%%%%%%%%%%%%%%%%%%%%%%%%%%%%%%%%%%%%%%%%%
\subsection{Stability}
\label{subsec:modeling-stability}
%%%%%%%%%%%%%%%%%%%%%%%%%%%%%%%%%%%%%%%%%%%%%%%%%%%%%%%%%%%%

A stable equilibrium resists sufficiently small perturbations.

An asymptotically stable equilibrium additionally attracts nearby
trajectories.

Stability analysis is central because a model may have an equilibrium that
is mathematically possible but physically unlikely to persist.

%%%%%%%%%%%%%%%%%%%%%%%%%%%%%%%%%%%%%%%%%%%%%%%%%%%%%%%%%%%%
\subsection{Thresholds}
\label{subsec:model-thresholds}
%%%%%%%%%%%%%%%%%%%%%%%%%%%%%%%%%%%%%%%%%%%%%%%%%%%%%%%%%%%%

Many systems exhibit threshold behavior.

A parameter may cross a critical value and change qualitative system
behavior.

For example:

\[
\boxed{
R_0<1
}
\]

may correspond to disease decline, while

\[
\boxed{
R_0>1
}
\]

may permit epidemic growth in a simple epidemic setting.

Thresholds often provide interpretable policy criteria.

%%%%%%%%%%%%%%%%%%%%%%%%%%%%%%%%%%%%%%%%%%%%%%%%%%%%%%%%%%%%
\subsection{Tipping Points}
\label{subsec:model-tipping-points}
%%%%%%%%%%%%%%%%%%%%%%%%%%%%%%%%%%%%%%%%%%%%%%%%%%%%%%%%%%%%

A tipping point is a threshold near which a small change in conditions can
produce a large qualitative change in the system.

Mathematically, tipping behavior may be associated with:

\[
\boxed{
\text{bifurcations},
}
\]

\[
\boxed{
\text{multiple equilibria},
}
\]

or

\[
\boxed{
\text{loss of stability}.
}
\]

%%%%%%%%%%%%%%%%%%%%%%%%%%%%%%%%%%%%%%%%%%%%%%%%%%%%%%%%%%%%
\subsection{Delay Models}
\label{subsec:model-delay-effects}
%%%%%%%%%%%%%%%%%%%%%%%%%%%%%%%%%%%%%%%%%%%%%%%%%%%%%%%%%%%%

Some systems respond to past rather than only current states.

A delay model may take form

\[
\boxed{
\frac{dx}{dt}
=
F
\left(
x(t),
x(t-\tau)
\right),
}
\]

where

\[
\tau>0
\]

is a delay.

Delays can generate oscillations and destabilize otherwise stable systems.

%%%%%%%%%%%%%%%%%%%%%%%%%%%%%%%%%%%%%%%%%%%%%%%%%%%%%%%%%%%%
\subsection{Spatial Models}
\label{subsec:model-spatial-models}
%%%%%%%%%%%%%%%%%%%%%%%%%%%%%%%%%%%%%%%%%%%%%%%%%%%%%%%%%%%%

If system behavior varies over location, the state may depend on both space
and time:

\[
\boxed{
u=u(\mathbf{x},t).
}
\]

Partial differential equations, spatial statistics, cellular models, and
networks can all represent spatial structure.

%%%%%%%%%%%%%%%%%%%%%%%%%%%%%%%%%%%%%%%%%%%%%%%%%%%%%%%%%%%%
\subsection{Diffusion Model}
\label{subsec:model-diffusion}
%%%%%%%%%%%%%%%%%%%%%%%%%%%%%%%%%%%%%%%%%%%%%%%%%%%%%%%%%%%%

A basic diffusion equation is

\[
\boxed{
\frac{\partial u}{\partial t}
=
D\nabla^2u,
}
\]

where

\[
D>0
\]

is a diffusion coefficient.

The model describes spreading caused by spatial concentration differences.

%%%%%%%%%%%%%%%%%%%%%%%%%%%%%%%%%%%%%%%%%%%%%%%%%%%%%%%%%%%%
\subsection{Advection-Diffusion Model}
\label{subsec:model-advection-diffusion}
%%%%%%%%%%%%%%%%%%%%%%%%%%%%%%%%%%%%%%%%%%%%%%%%%%%%%%%%%%%%

Transport by both flow and diffusion can be modeled by

\[
\boxed{
\frac{\partial u}{\partial t}
+
\mathbf{v}\cdot\nabla u
=
D\nabla^2u.
}
\]

Here:

\[
\mathbf{v}
=
\text{velocity field},
\]

and

\[
D
=
\text{diffusion coefficient}.
\]

Such equations appear in fluid, environmental, and transport modeling.

%%%%%%%%%%%%%%%%%%%%%%%%%%%%%%%%%%%%%%%%%%%%%%%%%%%%%%%%%%%%
\subsection{Network Models}
\label{subsec:model-network-models}
%%%%%%%%%%%%%%%%%%%%%%%%%%%%%%%%%%%%%%%%%%%%%%%%%%%%%%%%%%%%

Some systems are better described through relationships than physical
distance.

A network is represented by a graph

\[
\boxed{
G=(V,E).
}
\]

Vertices represent entities.

Edges represent interactions.

Examples include:

\[
\boxed{
\text{transportation networks},
}
\]

\[
\boxed{
\text{social networks},
}
\]

\[
\boxed{
\text{power grids},
}
\]

and

\[
\boxed{
\text{biological networks}.
}
\]

%%%%%%%%%%%%%%%%%%%%%%%%%%%%%%%%%%%%%%%%%%%%%%%%%%%%%%%%%%%%
\subsection{Dynamics on Networks}
\label{subsec:model-network-dynamics}
%%%%%%%%%%%%%%%%%%%%%%%%%%%%%%%%%%%%%%%%%%%%%%%%%%%%%%%%%%%%

A state

\[
x_i(t)
\]

may be assigned to each node.

Interactions can be modeled through adjacency matrix \(A\):

\[
\boxed{
\frac{dx_i}{dt}
=
F_i
\left(
x_i,
\sum_jA_{ij}x_j
\right).
}
\]

Network topology can strongly influence system dynamics.

%%%%%%%%%%%%%%%%%%%%%%%%%%%%%%%%%%%%%%%%%%%%%%%%%%%%%%%%%%%%
\subsection{Agent-Based Models}
\label{subsec:agent-based-models}
%%%%%%%%%%%%%%%%%%%%%%%%%%%%%%%%%%%%%%%%%%%%%%%%%%%%%%%%%%%%

Agent-based models represent individual entities explicitly.

Each agent follows behavioral rules.

System-level patterns then emerge from many local interactions.

Agents may represent:

\[
\boxed{
\text{people},
}
\]

\[
\boxed{
\text{animals},
}
\]

\[
\boxed{
\text{firms},
}
\]

or

\[
\boxed{
\text{vehicles}.
}
\]

%%%%%%%%%%%%%%%%%%%%%%%%%%%%%%%%%%%%%%%%%%%%%%%%%%%%%%%%%%%%
\subsection{Emergent Behavior}
\label{subsec:emergent-behavior}
%%%%%%%%%%%%%%%%%%%%%%%%%%%%%%%%%%%%%%%%%%%%%%%%%%%%%%%%%%%%

Emergent behavior is system-level behavior not explicitly prescribed at the
individual level.

Examples include:

\[
\boxed{
\text{traffic congestion},
}
\]

\[
\boxed{
\text{flocking},
}
\]

\[
\boxed{
\text{market patterns},
}
\]

and

\[
\boxed{
\text{epidemic waves}.
}
\]

Agent-based models are particularly useful for studying emergence.

%%%%%%%%%%%%%%%%%%%%%%%%%%%%%%%%%%%%%%%%%%%%%%%%%%%%%%%%%%%%
\subsection{Multiscale Models}
\label{subsec:multiscale-models}
%%%%%%%%%%%%%%%%%%%%%%%%%%%%%%%%%%%%%%%%%%%%%%%%%%%%%%%%%%%%

A multiscale model connects phenomena occurring at different scales.

For example:

\[
\boxed{
\text{molecular}
\rightarrow
\text{cellular}
\rightarrow
\text{tissue}
\rightarrow
\text{organ}.
}
\]

Different scales may require different mathematical models.

Coupling them is often one of the most difficult aspects of applied
mathematics.

%%%%%%%%%%%%%%%%%%%%%%%%%%%%%%%%%%%%%%%%%%%%%%%%%%%%%%%%%%%%
\subsection{Model Reduction}
\label{subsec:model-reduction}
%%%%%%%%%%%%%%%%%%%%%%%%%%%%%%%%%%%%%%%%%%%%%%%%%%%%%%%%%%%%

A complex model may contain more detail than needed.

Model reduction seeks a simpler system preserving important behavior.

For example,

\[
\boxed{
\dot{\mathbf{x}}
=
A\mathbf{x}
}
\]

with a large state dimension may be approximated in a lower-dimensional
subspace.

Reduction improves interpretation and computational efficiency.

%%%%%%%%%%%%%%%%%%%%%%%%%%%%%%%%%%%%%%%%%%%%%%%%%%%%%%%%%%%%
\subsection{Surrogate Models}
\label{subsec:surrogate-models}
%%%%%%%%%%%%%%%%%%%%%%%%%%%%%%%%%%%%%%%%%%%%%%%%%%%%%%%%%%%%

A surrogate model approximates a computationally expensive model.

Suppose

\[
F(\theta)
\]

requires hours of simulation.

One may construct an approximation

\[
\boxed{
\widehat{F}(\theta)
}
\]

that is much faster to evaluate.

Surrogates are useful for optimization, uncertainty analysis, and repeated
prediction.

%%%%%%%%%%%%%%%%%%%%%%%%%%%%%%%%%%%%%%%%%%%%%%%%%%%%%%%%%%%%
\subsection{Data-Driven Models}
\label{subsec:data-driven-models}
%%%%%%%%%%%%%%%%%%%%%%%%%%%%%%%%%%%%%%%%%%%%%%%%%%%%%%%%%%%%

A data-driven model learns relationships primarily from observed data.

Examples include:

\[
\boxed{
\text{regression},
}
\]

\[
\boxed{
\text{classification},
}
\]

\[
\boxed{
\text{time-series models},
}
\]

and

\[
\boxed{
\text{machine-learning models}.
}
\]

Their predictive ability depends strongly on the quality and relevance of
the data.

%%%%%%%%%%%%%%%%%%%%%%%%%%%%%%%%%%%%%%%%%%%%%%%%%%%%%%%%%%%%
\subsection{Mechanistic Versus Data-Driven Tradeoff}
\label{subsec:mechanistic-data-driven-tradeoff}
%%%%%%%%%%%%%%%%%%%%%%%%%%%%%%%%%%%%%%%%%%%%%%%%%%%%%%%%%%%%

Mechanistic models can offer:

\[
\boxed{
\text{interpretability}
}
\]

and

\[
\boxed{
\text{scientific structure}.
}
\]

Data-driven models may offer:

\[
\boxed{
\text{flexibility}
}
\]

and

\[
\boxed{
\text{predictive power}.
}
\]

Hybrid approaches attempt to combine both advantages.

%%%%%%%%%%%%%%%%%%%%%%%%%%%%%%%%%%%%%%%%%%%%%%%%%%%%%%%%%%%%
\subsection{Optimization Models}
\label{subsec:modeling-optimization-models}
%%%%%%%%%%%%%%%%%%%%%%%%%%%%%%%%%%%%%%%%%%%%%%%%%%%%%%%%%%%%

Many applied problems ask for the best decision.

A general optimization model is

\[
\boxed{
\min_{\mathbf{x}}
f(\mathbf{x})
}
\]

subject to

\[
\boxed{
g_i(\mathbf{x})\leq0
}
\]

and

\[
\boxed{
h_j(\mathbf{x})=0.
}
\]

The variables represent decisions.

The constraints represent feasibility.

%%%%%%%%%%%%%%%%%%%%%%%%%%%%%%%%%%%%%%%%%%%%%%%%%%%%%%%%%%%%
\subsection{Worked Example: Resource Allocation}
\label{subsec:worked-resource-allocation-model}
%%%%%%%%%%%%%%%%%%%%%%%%%%%%%%%%%%%%%%%%%%%%%%%%%%%%%%%%%%%%

\begin{workedexamplebox}{Allocating a Limited Resource}

Suppose products \(A\) and \(B\) produce profits

\[
5
\]

and

\[
4
\]

per unit.

They require resource amounts:

\[
2x+y\leq100.
\]

A basic optimization model is

\[
\max
\quad
5x+4y
\]

subject to

\[
2x+y\leq100,
\]

\[
x,y\geq0.
\]

\begin{answerbox}
\[
\boxed{
\max_{x,y\geq0}
\{
5x+4y:
2x+y\leq100
\}.
}
\]
\end{answerbox}

The model converts a verbal resource-allocation problem into mathematics.

\end{workedexamplebox}

%%%%%%%%%%%%%%%%%%%%%%%%%%%%%%%%%%%%%%%%%%%%%%%%%%%%%%%%%%%%
\subsection{Model Selection}
\label{subsec:model-selection}
%%%%%%%%%%%%%%%%%%%%%%%%%%%%%%%%%%%%%%%%%%%%%%%%%%%%%%%%%%%%

Several models may explain the same data.

Model selection balances:

\[
\boxed{
\text{fit},
}
\]

\[
\boxed{
\text{complexity},
}
\]

\[
\boxed{
\text{interpretability},
}
\]

and

\[
\boxed{
\text{predictive performance}.
}
\]

There is rarely one metric sufficient for every modeling purpose.

%%%%%%%%%%%%%%%%%%%%%%%%%%%%%%%%%%%%%%%%%%%%%%%%%%%%%%%%%%%%
\subsection{Residual Analysis}
\label{subsec:model-residual-analysis}
%%%%%%%%%%%%%%%%%%%%%%%%%%%%%%%%%%%%%%%%%%%%%%%%%%%%%%%%%%%%

Residuals should be examined for systematic patterns.

If residuals show:

\[
\boxed{
\text{curvature},
}
\]

\[
\boxed{
\text{changing variance},
}
\]

\[
\boxed{
\text{serial dependence},
}
\]

or

\[
\boxed{
\text{spatial structure},
}
\]

then important model assumptions may be violated.

%%%%%%%%%%%%%%%%%%%%%%%%%%%%%%%%%%%%%%%%%%%%%%%%%%%%%%%%%%%%
\subsection{Model Error}
\label{subsec:model-error}
%%%%%%%%%%%%%%%%%%%%%%%%%%%%%%%%%%%%%%%%%%%%%%%%%%%%%%%%%%%%

Even with perfect parameter estimates, the model equations themselves may
be imperfect.

A conceptual decomposition is

\[
\boxed{
\text{observation}
=
\text{model prediction}
+
\text{measurement noise}
+
\text{model discrepancy}.
}
\]

Model discrepancy reflects missing or incorrect structure.

%%%%%%%%%%%%%%%%%%%%%%%%%%%%%%%%%%%%%%%%%%%%%%%%%%%%%%%%%%%%
\subsection{Robustness}
\label{subsec:model-robustness}
%%%%%%%%%%%%%%%%%%%%%%%%%%%%%%%%%%%%%%%%%%%%%%%%%%%%%%%%%%%%

A model is robust when its important conclusions do not change dramatically
under reasonable perturbations of assumptions or parameters.

Robustness analysis may vary:

\[
\boxed{
\text{parameters},
}
\]

\[
\boxed{
\text{initial conditions},
}
\]

\[
\boxed{
\text{data},
}
\]

or

\[
\boxed{
\text{model structure}.
}
\]

%%%%%%%%%%%%%%%%%%%%%%%%%%%%%%%%%%%%%%%%%%%%%%%%%%%%%%%%%%%%
\subsection{Model Comparison}
\label{subsec:model-comparison}
%%%%%%%%%%%%%%%%%%%%%%%%%%%%%%%%%%%%%%%%%%%%%%%%%%%%%%%%%%%%

Competing models may be compared using:

\[
\boxed{
\text{prediction error},
}
\]

\[
\boxed{
\text{likelihood},
}
\]

\[
\boxed{
\text{information criteria},
}
\]

\[
\boxed{
\text{cross-validation},
}
\]

or

\[
\boxed{
\text{scientific interpretability}.
}
\]

The best criterion depends on the goal.

%%%%%%%%%%%%%%%%%%%%%%%%%%%%%%%%%%%%%%%%%%%%%%%%%%%%%%%%%%%%
\subsection{Prediction Versus Explanation}
\label{subsec:model-prediction-explanation}
%%%%%%%%%%%%%%%%%%%%%%%%%%%%%%%%%%%%%%%%%%%%%%%%%%%%%%%%%%%%

A model can predict well without revealing the correct mechanism.

A mechanistic model can explain important relationships while having lower
short-term predictive accuracy.

Thus:

\[
\boxed{
\text{prediction}
}
\]

and

\[
\boxed{
\text{explanation}
}
\]

are related but distinct modeling goals.

%%%%%%%%%%%%%%%%%%%%%%%%%%%%%%%%%%%%%%%%%%%%%%%%%%%%%%%%%%%%
\subsection{Correlation Versus Causation}
\label{subsec:model-correlation-causation}
%%%%%%%%%%%%%%%%%%%%%%%%%%%%%%%%%%%%%%%%%%%%%%%%%%%%%%%%%%%%

An observed relationship

\[
X\leftrightarrow Y
\]

does not by itself imply

\[
\boxed{
X
\text{ causes }
Y.
}
\]

Causal interpretation requires additional assumptions, experimental design,
or causal structure.

A predictive model should not automatically be interpreted causally.

%%%%%%%%%%%%%%%%%%%%%%%%%%%%%%%%%%%%%%%%%%%%%%%%%%%%%%%%%%%%
\subsection{Model Assumption Testing}
\label{subsec:model-assumption-testing}
%%%%%%%%%%%%%%%%%%%%%%%%%%%%%%%%%%%%%%%%%%%%%%%%%%%%%%%%%%%%

Assumptions should be checked where possible.

Examples include testing or diagnosing:

\[
\boxed{
\text{linearity},
}
\]

\[
\boxed{
\text{independence},
}
\]

\[
\boxed{
\text{constant variance},
}
\]

or

\[
\boxed{
\text{normality}.
}
\]

Some assumptions are scientific rather than statistical and require domain
knowledge.

%%%%%%%%%%%%%%%%%%%%%%%%%%%%%%%%%%%%%%%%%%%%%%%%%%%%%%%%%%%%
\subsection{Model Limitations}
\label{subsec:model-limitations}
%%%%%%%%%%%%%%%%%%%%%%%%%%%%%%%%%%%%%%%%%%%%%%%%%%%%%%%%%%%%

Every model has limitations.

These may include:

\[
\boxed{
\text{restricted parameter ranges},
}
\]

\[
\boxed{
\text{limited data},
}
\]

\[
\boxed{
\text{simplifying assumptions},
}
\]

\[
\boxed{
\text{unknown mechanisms},
}
\]

or

\[
\boxed{
\text{computational approximations}.
}
\]

A model should be used only within a justified domain.

%%%%%%%%%%%%%%%%%%%%%%%%%%%%%%%%%%%%%%%%%%%%%%%%%%%%%%%%%%%%
\subsection{Model Failure}
\label{subsec:model-failure}
%%%%%%%%%%%%%%%%%%%%%%%%%%%%%%%%%%%%%%%%%%%%%%%%%%%%%%%%%%%%

A model may fail because:

\[
\boxed{
\text{the assumptions are invalid},
}
\]

\[
\boxed{
\text{important variables are omitted},
}
\]

\[
\boxed{
\text{parameters have changed},
}
\]

or

\[
\boxed{
\text{the system entered a new regime}.
}
\]

Failure can be scientifically useful because it reveals what the model
does not capture.

%%%%%%%%%%%%%%%%%%%%%%%%%%%%%%%%%%%%%%%%%%%%%%%%%%%%%%%%%%%%
\subsection{Ethics of Mathematical Models}
\label{subsec:modeling-ethics}
%%%%%%%%%%%%%%%%%%%%%%%%%%%%%%%%%%%%%%%%%%%%%%%%%%%%%%%%%%%%

Models can influence decisions about:

\[
\boxed{
\text{health},
}
\]

\[
\boxed{
\text{finance},
}
\]

\[
\boxed{
\text{transportation},
}
\]

\[
\boxed{
\text{education},
}
\]

and

\[
\boxed{
\text{public policy}.
}
\]

Modelers should therefore communicate assumptions, uncertainty, limitations,
and potential consequences clearly.

%%%%%%%%%%%%%%%%%%%%%%%%%%%%%%%%%%%%%%%%%%%%%%%%%%%%%%%%%%%%
\subsection{Bias in Models}
\label{subsec:model-bias}
%%%%%%%%%%%%%%%%%%%%%%%%%%%%%%%%%%%%%%%%%%%%%%%%%%%%%%%%%%%%

Bias may enter through:

\[
\boxed{
\text{data collection},
}
\]

\[
\boxed{
\text{variable selection},
}
\]

\[
\boxed{
\text{model assumptions},
}
\]

or

\[
\boxed{
\text{evaluation criteria}.
}
\]

Mathematical sophistication does not automatically remove bias from a
modeling process.

%%%%%%%%%%%%%%%%%%%%%%%%%%%%%%%%%%%%%%%%%%%%%%%%%%%%%%%%%%%%
\subsection{Communicating Model Results}
\label{subsec:communicating-model-results}
%%%%%%%%%%%%%%%%%%%%%%%%%%%%%%%%%%%%%%%%%%%%%%%%%%%%%%%%%%%%

A modeling report should communicate:

\begin{itemize}
    \item the question;
    \item the model structure;
    \item definitions of variables and parameters;
    \item assumptions;
    \item data sources;
    \item parameter-estimation methods;
    \item uncertainty;
    \item validation results;
    \item limitations;
    \item interpretation of outputs.
\end{itemize}

A numerical result without context can be misleading.

%%%%%%%%%%%%%%%%%%%%%%%%%%%%%%%%%%%%%%%%%%%%%%%%%%%%%%%%%%%%
\subsection{Modeling Diagrams}
\label{subsec:modeling-diagrams}
%%%%%%%%%%%%%%%%%%%%%%%%%%%%%%%%%%%%%%%%%%%%%%%%%%%%%%%%%%%%

A diagram can clarify model structure before equations are written.

For example, a basic flow model is

\[
\boxed{
\text{Input}
\longrightarrow
\text{State}
\longrightarrow
\text{Output}.
}
\]

A feedback model may be written

\[
\boxed{
\text{State}
\longrightarrow
\text{Output}
\longrightarrow
\text{Control}
\longrightarrow
\text{State}.
}
\]

Such diagrams help reveal missing variables and feedback loops.

%%%%%%%%%%%%%%%%%%%%%%%%%%%%%%%%%%%%%%%%%%%%%%%%%%%%%%%%%%%%
\subsection{Modeling With Data}
\label{subsec:modeling-with-data}
%%%%%%%%%%%%%%%%%%%%%%%%%%%%%%%%%%%%%%%%%%%%%%%%%%%%%%%%%%%%

Data may be used to:

\[
\boxed{
\text{estimate parameters},
}
\]

\[
\boxed{
\text{choose models},
}
\]

\[
\boxed{
\text{validate predictions},
}
\]

and

\[
\boxed{
\text{identify missing mechanisms}.
}
\]

The relationship between model and data is therefore iterative.

%%%%%%%%%%%%%%%%%%%%%%%%%%%%%%%%%%%%%%%%%%%%%%%%%%%%%%%%%%%%
\subsection{The Modeling Loop}
\label{subsec:modeling-loop}
%%%%%%%%%%%%%%%%%%%%%%%%%%%%%%%%%%%%%%%%%%%%%%%%%%%%%%%%%%%%

A useful modeling loop is:

\[
\boxed{
\text{observe}
}
\]

\[
\downarrow
\]

\[
\boxed{
\text{model}
}
\]

\[
\downarrow
\]

\[
\boxed{
\text{predict}
}
\]

\[
\downarrow
\]

\[
\boxed{
\text{test}
}
\]

\[
\downarrow
\]

\[
\boxed{
\text{revise}.
}
\]

Scientific modeling is therefore a process of repeated confrontation between
mathematical structure and evidence.

%%%%%%%%%%%%%%%%%%%%%%%%%%%%%%%%%%%%%%%%%%%%%%%%%%%%%%%%%%%%
\subsection{A Complete Worked Modeling Example}
\label{subsec:complete-worked-modeling-example}
%%%%%%%%%%%%%%%%%%%%%%%%%%%%%%%%%%%%%%%%%%%%%%%%%%%%%%%%%%%%

\begin{workedexamplebox}{Cooling of an Object}

Suppose an object with temperature \(T(t)\) is placed in a room with
constant ambient temperature \(T_a\).

Assume the rate of temperature change is proportional to the difference
between object and room temperature:

\[
\frac{dT}{dt}
=
-k(T-T_a),
\qquad
k>0.
\]

\textbf{Variables and parameters}

\[
T(t)
=
\text{object temperature},
\]

\[
T_a
=
\text{ambient temperature},
\]

\[
k
=
\text{cooling constant}.
\]

\textbf{Solve the model}

Rewrite:

\[
\frac{d}{dt}(T-T_a)
=
-k(T-T_a).
\]

Therefore,

\[
T-T_a
=
Ce^{-kt}.
\]

If

\[
T(0)=T_0,
\]

then

\[
C=T_0-T_a.
\]

Hence

\[
\boxed{
T(t)
=
T_a
+
(T_0-T_a)e^{-kt}.
}
\]

\textbf{Interpretation}

As

\[
t\to\infty,
\]

\[
e^{-kt}\to0,
\]

so

\[
\begin{answerbox}
\[
\boxed{
T(t)\to T_a.
}
\]
\end{answerbox}

The model predicts exponential approach to ambient temperature.

\end{workedexamplebox}

%%%%%%%%%%%%%%%%%%%%%%%%%%%%%%%%%%%%%%%%%%%%%%%%%%%%%%%%%%%%
\subsection{A Complete Data-Calibration Example}
\label{subsec:complete-calibration-example}
%%%%%%%%%%%%%%%%%%%%%%%%%%%%%%%%%%%%%%%%%%%%%%%%%%%%%%%%%%%%

\begin{workedexamplebox}{Estimating a Cooling Constant}

Suppose Newton cooling gives

\[
T(t)-T_a
=
(T_0-T_a)e^{-kt}.
\]

Assume

\[
T_a=20,
\]

\[
T_0=80,
\]

and after \(10\) minutes,

\[
T(10)=50.
\]

Then

\[
50-20
=
(80-20)e^{-10k}.
\]

Thus

\[
30
=
60e^{-10k}.
\]

Divide by \(60\):

\[
\frac12
=
e^{-10k}.
\]

Taking logarithms,

\[
-\ln2
=
-10k.
\]

Therefore,

\begin{answerbox}
\[
\boxed{
k
=
\frac{\ln2}{10}.
}
\]
\end{answerbox}

The observed temperature has calibrated the model parameter.

\end{workedexamplebox}

%%%%%%%%%%%%%%%%%%%%%%%%%%%%%%%%%%%%%%%%%%%%%%%%%%%%%%%%%%%%
\subsection{Modeling Hierarchy}
\label{subsec:modeling-hierarchy}
%%%%%%%%%%%%%%%%%%%%%%%%%%%%%%%%%%%%%%%%%%%%%%%%%%%%%%%%%%%%

A useful hierarchy of model complexity is:

\[
\boxed{
\text{algebraic}
}
\]

\[
\longrightarrow
\]

\[
\boxed{
\text{difference or differential equation}
}
\]

\[
\longrightarrow
\]

\[
\boxed{
\text{stochastic}
}
\]

\[
\longrightarrow
\]

\[
\boxed{
\text{spatial or network}
}
\]

\[
\longrightarrow
\]

\[
\boxed{
\text{multiscale or agent-based}.
}
\]

One should move to a more complicated model only when the added structure
serves the modeling objective.

%%%%%%%%%%%%%%%%%%%%%%%%%%%%%%%%%%%%%%%%%%%%%%%%%%%%%%%%%%%%
\subsection{Modeling Principles}
\label{subsec:modeling-principles}
%%%%%%%%%%%%%%%%%%%%%%%%%%%%%%%%%%%%%%%%%%%%%%%%%%%%%%%%%%%%

Useful general principles include:

\begin{enumerate}
    \item begin with the question, not the equation;
    \item define variables clearly;
    \item state assumptions explicitly;
    \item use the simplest model adequate for the purpose;
    \item preserve units and dimensions;
    \item separate parameters from state variables;
    \item calibrate with appropriate data;
    \item validate on information not used for fitting when possible;
    \item quantify uncertainty;
    \item examine sensitivity;
    \item test robustness;
    \item avoid unjustified extrapolation;
    \item distinguish prediction from explanation;
    \item distinguish association from causation;
    \item document limitations;
    \item revise the model when evidence requires it.
\end{enumerate}

%%%%%%%%%%%%%%%%%%%%%%%%%%%%%%%%%%%%%%%%%%%%%%%%%%%%%%%%%%%%
\subsection{Common Errors}
\label{subsec:mathematical-modeling-common-errors}
%%%%%%%%%%%%%%%%%%%%%%%%%%%%%%%%%%%%%%%%%%%%%%%%%%%%%%%%%%%%

\subsubsection{Treating the Model as Reality}

A mathematical model is a representation of selected features.

It is not identical to the system itself.

%%%%%%%%%%%%%%%%%%%%%%%%%%%%%%%%%%%%%%%%%%%%%%%%%%%%%%%%%%%%
\subsubsection{Choosing Equations Before Defining the Question}

A familiar equation should not be selected merely because it is convenient.

The model form should follow from the purpose and assumptions.

%%%%%%%%%%%%%%%%%%%%%%%%%%%%%%%%%%%%%%%%%%%%%%%%%%%%%%%%%%%%
\subsubsection{Failing to State Assumptions}

Hidden assumptions make it difficult to judge when a model is valid.

Important assumptions should be written explicitly.

%%%%%%%%%%%%%%%%%%%%%%%%%%%%%%%%%%%%%%%%%%%%%%%%%%%%%%%%%%%%
\subsubsection{Using Too Many Variables}

More detail is not automatically better.

Unnecessary variables can increase:

\[
\boxed{
\text{parameter uncertainty},
}
\]

\[
\boxed{
\text{computational cost},
}
\]

and

\[
\boxed{
\text{identifiability problems}.
}
\]

%%%%%%%%%%%%%%%%%%%%%%%%%%%%%%%%%%%%%%%%%%%%%%%%%%%%%%%%%%%%
\subsubsection{Using Too Few Variables}

A model can also be too simple.

If an omitted variable controls important behavior, the model may fail even
with excellent parameter fitting.

%%%%%%%%%%%%%%%%%%%%%%%%%%%%%%%%%%%%%%%%%%%%%%%%%%%%%%%%%%%%
\subsubsection{Ignoring Units}

An equation that is dimensionally inconsistent cannot correctly represent a
physical relationship.

Always verify that terms have compatible dimensions.

%%%%%%%%%%%%%%%%%%%%%%%%%%%%%%%%%%%%%%%%%%%%%%%%%%%%%%%%%%%%
\subsubsection{Confusing Parameters and Variables}

State variables evolve according to the model.

Parameters are usually treated as fixed during one model run.

The distinction should be explicit.

%%%%%%%%%%%%%%%%%%%%%%%%%%%%%%%%%%%%%%%%%%%%%%%%%%%%%%%%%%%%
\subsubsection{Confusing Calibration With Validation}

Calibration fits the model.

Validation evaluates its performance for the intended purpose.

Using the same data for both without care can overstate model quality.

%%%%%%%%%%%%%%%%%%%%%%%%%%%%%%%%%%%%%%%%%%%%%%%%%%%%%%%%%%%%
\subsubsection{Confusing Verification With Validation}

Verification asks whether the model was implemented correctly.

Validation asks whether the model adequately represents reality for its
purpose.

%%%%%%%%%%%%%%%%%%%%%%%%%%%%%%%%%%%%%%%%%%%%%%%%%%%%%%%%%%%%
\subsubsection{Assuming Good Fit Proves the Mechanism}

Different models can fit the same data.

Statistical agreement does not by itself establish a causal mechanism.

%%%%%%%%%%%%%%%%%%%%%%%%%%%%%%%%%%%%%%%%%%%%%%%%%%%%%%%%%%%%
\subsubsection{Ignoring Identifiability}

An optimization algorithm may return parameter estimates even when the data
cannot uniquely determine those parameters.

Numerical output does not guarantee identifiability.

%%%%%%%%%%%%%%%%%%%%%%%%%%%%%%%%%%%%%%%%%%%%%%%%%%%%%%%%%%%%
\subsubsection{Ignoring Parameter Correlation}

Two parameters may compensate for one another.

This can create broad valleys in an objective function and large uncertainty
even if the fitted output looks accurate.

%%%%%%%%%%%%%%%%%%%%%%%%%%%%%%%%%%%%%%%%%%%%%%%%%%%%%%%%%%%%
\subsubsection{Ignoring Sensitivity}

If a small change in one parameter produces a huge output change, predictions
may be fragile.

Sensitivity should be investigated before strong conclusions are drawn.

%%%%%%%%%%%%%%%%%%%%%%%%%%%%%%%%%%%%%%%%%%%%%%%%%%%%%%%%%%%%
\subsubsection{Ignoring Structural Uncertainty}

Parameter uncertainty assumes the model structure is fixed.

The equations themselves may also be wrong or incomplete.

%%%%%%%%%%%%%%%%%%%%%%%%%%%%%%%%%%%%%%%%%%%%%%%%%%%%%%%%%%%%
\subsubsection{Extrapolating Too Far}

A model calibrated on one domain should not automatically be trusted far
outside that domain.

Extrapolation requires additional justification.

%%%%%%%%%%%%%%%%%%%%%%%%%%%%%%%%%%%%%%%%%%%%%%%%%%%%%%%%%%%%
\subsubsection{Confusing Deterministic and Certain}

A deterministic model can still have uncertain parameters or uncertain
initial conditions.

Deterministic does not mean perfectly known.

%%%%%%%%%%%%%%%%%%%%%%%%%%%%%%%%%%%%%%%%%%%%%%%%%%%%%%%%%%%%
\subsubsection{Confusing Stochastic With Unstructured}

A stochastic model still has mathematical structure.

Probability distributions and dependencies are part of the model.

%%%%%%%%%%%%%%%%%%%%%%%%%%%%%%%%%%%%%%%%%%%%%%%%%%%%%%%%%%%%
\subsubsection{Assuming Simulation Is the Model}

The model consists of mathematical assumptions and equations.

A simulation is one computational method used to investigate the model.

%%%%%%%%%%%%%%%%%%%%%%%%%%%%%%%%%%%%%%%%%%%%%%%%%%%%%%%%%%%%
\subsubsection{Assuming a Simulation Run Is a Proof}

A numerical trajectory provides evidence about one computational case.

It does not automatically establish a theorem about every parameter value or
initial condition.

%%%%%%%%%%%%%%%%%%%%%%%%%%%%%%%%%%%%%%%%%%%%%%%%%%%%%%%%%%%%
\subsubsection{Ignoring Numerical Error}

A simulation approximates the mathematical model.

Discretization, rounding, and solver tolerances can introduce additional
error.

%%%%%%%%%%%%%%%%%%%%%%%%%%%%%%%%%%%%%%%%%%%%%%%%%%%%%%%%%%%%
\subsubsection{Confusing Prediction With Causation}

A model may predict \(Y\) from \(X\) accurately without showing that changing
\(X\) causes \(Y\) to change.

Causal claims require stronger assumptions.

%%%%%%%%%%%%%%%%%%%%%%%%%%%%%%%%%%%%%%%%%%%%%%%%%%%%%%%%%%%%
\subsubsection{Overinterpreting Precision}

A model output such as

\[
17.384729
\]

does not imply the real system is known to six decimal places.

Reported precision should reflect data and model uncertainty.

%%%%%%%%%%%%%%%%%%%%%%%%%%%%%%%%%%%%%%%%%%%%%%%%%%%%%%%%%%%%
\subsubsection{Ignoring the Modeling Purpose}

A model can be excellent for one purpose and poor for another.

Validation must be tied to intended use.

%%%%%%%%%%%%%%%%%%%%%%%%%%%%%%%%%%%%%%%%%%%%%%%%%%%%%%%%%%%%
\subsection{Applications}
\label{subsec:mathematical-modeling-applications}
%%%%%%%%%%%%%%%%%%%%%%%%%%%%%%%%%%%%%%%%%%%%%%%%%%%%%%%%%%%%

\begin{applicationbox}

\textbf{Biology.}
Population growth, ecosystems, disease transmission, genetics, and cellular
processes can be represented using differential equations, stochastic
models, and networks.

\medskip

\textbf{Physics.}
Mechanics, transport, waves, fields, thermodynamics, and quantum systems rely
heavily on mathematical modeling.

\medskip

\textbf{Engineering.}
Models support design, control, reliability analysis, optimization, and
simulation.

\medskip

\textbf{Environmental Science.}
Pollutant transport, climate processes, hydrology, and ecosystem dynamics
require multiscale spatial and temporal models.

\medskip

\textbf{Economics.}
Demand, production, growth, equilibrium, strategic interaction, and dynamic
decision-making can be formulated mathematically.

\medskip

\textbf{Finance.}
Asset prices, portfolio risk, derivatives, interest rates, and default can
be modeled probabilistically and dynamically.

\medskip

\textbf{Operations Research.}
Scheduling, transportation, inventory, routing, and resource allocation are
modeled through optimization.

\medskip

\textbf{Data Science.}
Regression, statistical learning, simulation, calibration, and uncertainty
quantification provide data-driven mathematical models.

\medskip

\textbf{Networks.}
Transportation, communication, biological, social, and infrastructure
systems can be modeled as interacting nodes and edges.

\end{applicationbox}

%%%%%%%%%%%%%%%%%%%%%%%%%%%%%%%%%%%%%%%%%%%%%%%%%%%%%%%%%%%%
\subsection{Exercises}
\label{subsec:mathematical-modeling-exercises}
%%%%%%%%%%%%%%%%%%%%%%%%%%%%%%%%%%%%%%%%%%%%%%%%%%%%%%%%%%%%

\begin{exercise}[1]
Define a mathematical model.
\end{exercise}

\begin{exercise}[1]
Explain why a model is not identical to the system it represents.
\end{exercise}

\begin{exercise}[1]
Define a state variable.
\end{exercise}

\begin{exercise}[1]
Define a model parameter.
\end{exercise}

\begin{exercise}[1]
Distinguish endogenous and exogenous variables.
\end{exercise}

\begin{exercise}[1]
Define a mechanistic model.
\end{exercise}

\begin{exercise}[1]
Define an empirical model.
\end{exercise}

\begin{exercise}[1]
Distinguish deterministic and stochastic models.
\end{exercise}

\begin{exercise}[1]
Distinguish static and dynamic models.
\end{exercise}

\begin{exercise}[1]
Distinguish discrete and continuous-time models.
\end{exercise}

\begin{exercise}[2]
Give three assumptions that might appear in a population model.
\end{exercise}

\begin{exercise}[2]
For

\[
P_{n+1}=1.03P_n,
\]

identify the state variable and parameter.
\end{exercise}

\begin{exercise}[2]
Solve

\[
P_{n+1}=1.03P_n,
\qquad
P_0=1000.
\]
\end{exercise}

\begin{exercise}[2]
Construct a continuous exponential-growth model from the statement:
``the population growth rate is proportional to population size.''
\end{exercise}

\begin{exercise}[2]
Solve

\[
\frac{dP}{dt}=0.05P,
\qquad
P(0)=200.
\]
\end{exercise}

\begin{exercise}[2]
Find the equilibria of

\[
\frac{dP}{dt}
=
rP
\left(
1-\frac{P}{K}
\right).
\]
\end{exercise}

\begin{exercise}[3]
Analyze the sign of logistic growth for:

\[
0<P<K,
\]

and

\[
P>K.
\]
\end{exercise}

\begin{exercise}[3]
Show that logistic growth is maximized at

\[
P=\frac{K}{2}.
\]
\end{exercise}

\begin{exercise}[2]
Write a general conservation balance equation.
\end{exercise}

\begin{exercise}[3]
Construct a salt-mixing model for a tank with unequal inflow and outflow
rates.
\end{exercise}

\begin{exercise}[3]
Construct a two-compartment transfer model.
\end{exercise}

\begin{exercise}[3]
Show conservation of total quantity in a closed two-compartment model.
\end{exercise}

\begin{exercise}[2]
Check the dimensional consistency of

\[
v=v_0+at.
\]
\end{exercise}

\begin{exercise}[2]
Check the dimensional consistency of

\[
x=x_0+v_0t+\frac12at^2.
\]
\end{exercise}

\begin{exercise}[3]
Explain why terms with incompatible units cannot be added.
\end{exercise}

\begin{exercise}[3]
Nondimensionalize

\[
\frac{dP}{dt}
=
rP
\left(
1-\frac{P}{K}
\right)
\]

using

\[
u=\frac{P}{K},
\qquad
\tau=rt.
\]
\end{exercise}

\begin{exercise}[3]
Explain why nondimensionalization can reduce the number of parameters.
\end{exercise}

\begin{exercise}[2]
For the power law

\[
Y=cX^\alpha,
\]

derive the corresponding linear log-log relationship.
\end{exercise}

\begin{exercise}[2]
Define a residual.
\end{exercise}

\begin{exercise}[3]
Derive the least-squares estimator for

\[
y=ax
\]

with no intercept.
\end{exercise}

\begin{exercise}[2]
Distinguish calibration and validation.
\end{exercise}

\begin{exercise}[2]
Distinguish verification and validation.
\end{exercise}

\begin{exercise}[3]
Give an example of a model that could be verified computationally but fail
validation.
\end{exercise}

\begin{exercise}[3]
Explain overfitting.
\end{exercise}

\begin{exercise}[3]
Explain underfitting.
\end{exercise}

\begin{exercise}[3]
Explain parsimony.
\end{exercise}

\begin{exercise}[2]
Distinguish interpolation and extrapolation.
\end{exercise}

\begin{exercise}[3]
Explain why extrapolation can be dangerous even with a small training error.
\end{exercise}

\begin{exercise}[2]
Define identifiability.
\end{exercise}

\begin{exercise}[3]
Distinguish structural and practical identifiability.
\end{exercise}

\begin{exercise}[3]
Analyze the identifiability of parameters \(a,b\) in

\[
y(t)=ab\,t.
\]
\end{exercise}

\begin{exercise}[2]
Define local sensitivity.
\end{exercise}

\begin{exercise}[3]
For

\[
y=\theta^2,
\]

compute

\[
\frac{dy}{d\theta}.
\]
\end{exercise}

\begin{exercise}[3]
For

\[
P(t)=P_0e^{rt},
\]

compute sensitivity with respect to \(P_0\).
\end{exercise}

\begin{exercise}[3]
Compute normalized sensitivity of

\[
P(t)=P_0e^{rt}
\]

with respect to \(r\).
\end{exercise}

\begin{exercise}[3]
Distinguish local and global sensitivity.
\end{exercise}

\begin{exercise}[2]
Define uncertainty quantification conceptually.
\end{exercise}

\begin{exercise}[3]
List four sources of uncertainty in a model.
\end{exercise}

\begin{exercise}[3]
Distinguish aleatoric and epistemic uncertainty.
\end{exercise}

\begin{exercise}[2]
Explain scenario analysis.
\end{exercise}

\begin{exercise}[2]
Define simulation.
\end{exercise}

\begin{exercise}[3]
Explain how Monte Carlo simulation estimates an expectation.
\end{exercise}

\begin{exercise}[2]
Define an equilibrium of

\[
\dot{x}=f(x).
\]
\end{exercise}

\begin{exercise}[3]
Find the equilibrium of

\[
\frac{dx}{dt}
=
-k(x-a).
\]
\end{exercise}

\begin{exercise}[3]
Solve

\[
\frac{dx}{dt}
=
-k(x-a),
\qquad
x(0)=x_0.
\]
\end{exercise}

\begin{exercise}[3]
Explain why this is a negative-feedback model.
\end{exercise}

\begin{exercise}[3]
Give an example of positive feedback.
\end{exercise}

\begin{exercise}[3]
Explain threshold behavior.
\end{exercise}

\begin{exercise}[3]
Explain a tipping point conceptually.
\end{exercise}

\begin{exercise}[2]
Write a generic delay differential equation.
\end{exercise}

\begin{exercise}[3]
Explain why delay can change stability.
\end{exercise}

\begin{exercise}[2]
Write the diffusion equation.
\end{exercise}

\begin{exercise}[2]
Explain the role of diffusion coefficient \(D\).
\end{exercise}

\begin{exercise}[3]
Write an advection-diffusion equation.
\end{exercise}

\begin{exercise}[2]
Define a network model.
\end{exercise}

\begin{exercise}[3]
Explain how an adjacency matrix encodes interactions.
\end{exercise}

\begin{exercise}[3]
Construct a simple linear network-diffusion model conceptually.
\end{exercise}

\begin{exercise}[2]
Define an agent-based model.
\end{exercise}

\begin{exercise}[3]
Explain emergent behavior.
\end{exercise}

\begin{exercise}[3]
Give an example of a system where agent-based modeling is preferable to a
single aggregate differential equation.
\end{exercise}

\begin{exercise}[2]
Define a multiscale model.
\end{exercise}

\begin{exercise}[3]
Give an example of two scales that must be coupled in a biological model.
\end{exercise}

\begin{exercise}[2]
Explain model reduction.
\end{exercise}

\begin{exercise}[2]
Explain a surrogate model.
\end{exercise}

\begin{exercise}[3]
Explain why surrogate models are useful in optimization.
\end{exercise}

\begin{exercise}[3]
Compare mechanistic and data-driven models.
\end{exercise}

\begin{exercise}[3]
Explain a hybrid mechanistic-data-driven model.
\end{exercise}

\begin{exercise}[2]
Formulate a two-variable linear resource-allocation problem.
\end{exercise}

\begin{exercise}[3]
Distinguish model fit from causal interpretation.
\end{exercise}

\begin{exercise}[3]
Explain why correlation does not imply causation.
\end{exercise}

\begin{exercise}[3]
List several patterns that may appear in residual plots when a model is
misspecified.
\end{exercise}

\begin{exercise}[3]
Explain model discrepancy.
\end{exercise}

\begin{exercise}[3]
Explain robustness.
\end{exercise}

\begin{exercise}[3]
Describe how you would test robustness of a parameterized model.
\end{exercise}

\begin{exercise}[3]
Explain why more decimal places in model output do not imply greater real
accuracy.
\end{exercise}

\begin{exercise}[3]
Explain why a model can be useful even if it is not literally true.
\end{exercise}

\begin{exercise}[3]
Describe the complete mathematical-modeling workflow.
\end{exercise}

\begin{exercise}[3]
Construct a model for a bank account with annual compound interest.
\end{exercise}

\begin{exercise}[3]
Construct a continuous model for radioactive decay.
\end{exercise}

\begin{exercise}[3]
Construct a model for a reservoir receiving a constant inflow and losing
water at a rate proportional to current volume.
\end{exercise}

\begin{exercise}[3]
Construct a model for temperature approaching ambient temperature.
\end{exercise}

\begin{exercise}[3]
Construct a logistic population model and explain each parameter.
\end{exercise}

\begin{exercise}[3]
Construct a simple supply-demand equilibrium model.
\end{exercise}

\begin{exercise}[3]
Construct a basic linear regression model and identify which terms represent
systematic behavior and noise.
\end{exercise}

\begin{exercise}[3]
Construct a two-state compartment model with external input and removal.
\end{exercise}

\begin{exercise}[4]
Develop a full parameter-estimation procedure for a nonlinear ODE model.
\end{exercise}

\begin{exercise}[4]
Study structural identifiability of a two-parameter dynamical model.
\end{exercise}

\begin{exercise}[4]
Perform a local sensitivity analysis using sensitivity equations.
\end{exercise}

\begin{exercise}[4]
Design a Monte Carlo uncertainty-quantification experiment.
\end{exercise}

\begin{exercise}[4]
Compare deterministic and stochastic versions of the same population model.
\end{exercise}

\begin{exercise}[4]
Develop and nondimensionalize a predator-prey model.
\end{exercise}

\begin{exercise}[4]
Construct and analyze a compartmental epidemic model.
\end{exercise}

\begin{exercise}[4]
Develop a network diffusion model using a graph Laplacian.
\end{exercise}

\begin{exercise}[4]
Develop a simple agent-based epidemic model.
\end{exercise}

\begin{exercise}[4]
Compare an agent-based epidemic model with an aggregate compartment model.
\end{exercise}

\begin{exercise}[4]
Construct a surrogate model for an expensive simulation.
\end{exercise}

\begin{exercise}[4]
Develop a model comparison study using held-out data.
\end{exercise}

\begin{exercise}[4]
Analyze a model under several competing structural assumptions.
\end{exercise}

\begin{exercise}[4]
Study how uncertainty in parameters propagates to uncertainty in model
outputs.
\end{exercise}

\begin{exercise}[4]
Construct a multiobjective modeling problem balancing accuracy, cost, and
risk.
\end{exercise}

%%%%%%%%%%%%%%%%%%%%%%%%%%%%%%%%%%%%%%%%%%%%%%%%%%%%%%%%%%%%
\subsection{Conceptual Summary}
\label{subsec:mathematical-modeling-summary}
%%%%%%%%%%%%%%%%%%%%%%%%%%%%%%%%%%%%%%%%%%%%%%%%%%%%%%%%%%%%

Mathematical modeling translates selected features of a system into
mathematical structure.

A model typically contains:

\[
\boxed{
\text{variables}
}
\]

\[
\boxed{
+
}
\]

\[
\boxed{
\text{parameters}
}
\]

\[
\boxed{
+
}
\]

\[
\boxed{
\text{assumptions}
}
\]

\[
\boxed{
+
}
\]

\[
\boxed{
\text{equations or algorithms}.
}
\]

Models may be:

\[
\boxed{
\text{mechanistic or empirical},
}
\]

\[
\boxed{
\text{deterministic or stochastic},
}
\]

\[
\boxed{
\text{discrete or continuous},
}
\]

\[
\boxed{
\text{static or dynamic},
}
\]

and

\[
\boxed{
\text{linear or nonlinear}.
}
\]

Important modeling tools include:

\[
\boxed{
\text{conservation laws},
}
\]

\[
\boxed{
\text{dimensional analysis},
}
\]

\[
\boxed{
\text{parameter estimation},
}
\]

\[
\boxed{
\text{sensitivity analysis},
}
\]

\[
\boxed{
\text{uncertainty quantification},
}
\]

and

\[
\boxed{
\text{simulation}.
}
\]

A useful modeling cycle is:

\[
\boxed{
\text{question}
\rightarrow
\text{model}
\rightarrow
\text{analysis}
\rightarrow
\text{data}
\rightarrow
\text{revision}.
}
\]

A model should not be judged only by mathematical elegance.

It must be appropriate for its intended purpose.

The central question is therefore not:

\[
\boxed{
\text{Is the model perfectly true?}
}
\]

but rather:

\[
\boxed{
\text{Is the model sufficiently accurate, interpretable, and useful for the
question being asked?}
}
\]

%%%%%%%%%%%%%%%%%%%%%%%%%%%%%%%%%%%%%%%%%%%%%%%%%%%%%%%%%%%%
\subsection{Section Connections}
\label{subsec:mathematical-modeling-connections}
%%%%%%%%%%%%%%%%%%%%%%%%%%%%%%%%%%%%%%%%%%%%%%%%%%%%%%%%%%%%

\chapterconnection{
Mathematical Modeling begins Chapter~13 by connecting the abstract theories
developed throughout this text with real systems and applied questions.
Algebra provides symbolic relationships, calculus and differential
equations describe change, probability and statistics handle uncertainty
and data, optimization identifies preferred decisions, numerical mathematics
makes complicated models computationally tractable, and networks describe
interacting systems. The remainder of Chapter~13 develops specialized
modeling frameworks for biology, physics, continuum mechanics, fluid
dynamics, control, networks, finance, and economics. The next section,
Mathematical Biology, applies the modeling process to populations,
epidemiology, ecology, genetics, biochemical systems, and other biological
phenomena.
}

%%%%%%%%%%%%%%%%%%%%%%%%%%%%%%%%%%%%%%%%%%%%%%%%%%%%%%%%%%%%
% SECTION SOURCE TRACKING
%%%%%%%%%%%%%%%%%%%%%%%%%%%%%%%%%%%%%%%%%%%%%%%%%%%%%%%%%%%%

%------------------------------------------------------------
% 13.1 Mathematical Modeling
%------------------------------------------------------------
%
% Source basis:
%   - Standard applied mathematics and mathematical modeling principles.
%   - Standard differential-equation, statistical, optimization,
%     numerical, and scientific-modeling concepts.
%
% Used for:
%   - definition and purpose of mathematical models;
%   - abstraction;
%   - variables and parameters;
%   - states, inputs, outputs;
%   - endogenous and exogenous quantities;
%   - assumptions;
%   - mechanistic models;
%   - empirical models;
%   - hybrid models;
%   - deterministic and stochastic models;
%   - discrete and continuous models;
%   - static and dynamic models;
%   - linear and nonlinear models;
%   - modeling workflow;
%   - algebraic models;
%   - difference equations;
%   - differential equations;
%   - exponential growth;
%   - logistic growth;
%   - conservation laws;
%   - balance equations;
%   - compartment models;
%   - dimensional analysis;
%   - dimensional consistency;
%   - nondimensionalization;
%   - scaling laws;
%   - parameter estimation;
%   - residuals;
%   - calibration;
%   - validation;
%   - verification;
%   - training and testing;
%   - overfitting and underfitting;
%   - parsimony;
%   - interpolation and extrapolation;
%   - identifiability;
%   - structural and practical identifiability;
%   - local and global sensitivity;
%   - normalized sensitivities;
%   - uncertainty quantification;
%   - aleatoric and epistemic uncertainty;
%   - scenario analysis;
%   - simulation;
%   - Monte Carlo modeling;
%   - equilibria;
%   - feedback;
%   - stability;
%   - thresholds and tipping points;
%   - delays;
%   - spatial modeling;
%   - diffusion and advection;
%   - network models;
%   - agent-based models;
%   - emergent behavior;
%   - multiscale models;
%   - model reduction;
%   - surrogate models;
%   - data-driven models;
%   - optimization models;
%   - model selection;
%   - residual analysis;
%   - model discrepancy;
%   - robustness;
%   - prediction versus explanation;
%   - correlation versus causation;
%   - model limitations;
%   - modeling ethics;
%   - communication;
%   - applications and exercises.
%
% Notes:
%   - Mathematical Biology is developed next in Section 13.2.
%   - Mathematical Physics follows in Section 13.3.
%   - Continuum Mechanics follows in Section 13.4.
%   - Fluid Dynamics follows in Section 13.5.
%   - Control Theory follows in Section 13.6.
%   - Network Science follows in Section 13.7.
%   - Mathematical Finance follows in Section 13.8.
%   - Mathematical Economics follows in Section 13.9.
%
%%%%%%%%%%%%%%%%%%%%%%%%%%%%%%%%%%%%%%%%%%%%%%%%%%%%%%%%%%%%